% ===========================================================================
%  chapter.tex --- thesis-ready fragment, journal-article form.
%  Requires the macros of preamble.tex.
%
%  This is a rewritten, condensed version of the population chapter,
%  organised as a research article: one construction, three consequences,
%  and a scope section.  The birth--death--catastrophe process and the
%  comparison with budding are the spine; the flooding parameter L carries
%  the central biological result.
% ===========================================================================

\chapter[Burst-aware within-host dynamics]{A birth--death--catastrophe
extension of the standard within-host model, and what it says about
bursting}
\label{p:ch:one-cell-to-population}

\section{Introduction}
\label{p:sec:introduction}

The standard model of within-host virus dynamics describes an infected cell
with two numbers: it releases virions at a constant rate $p$ and is removed
at a constant rate $d_{\Icell}$
\cite{mclean1993balance,nowak1996population,perelson1996hiv}. Both are
fitted to viral load data and then reported as though they characterised the
pathogen. They are, however, phenomenological. Nothing in the standard model
says what the cell is doing between infection and release, and the constants
absorb whatever it is.

For one class of pathogen this costs nothing. A budding virus really does
experience release and death as separate, independent events, and
\cref{p:sec:budding} shows that the standard model is then an exact
description of the cell, not an approximation to it. For a lytic pathogen it
is wrong in mechanism. Such a pathogen releases nothing while it replicates
inside the cell and then releases everything at once, in the event that ends
the cell's productive life. Release and death are not independent; they are
the same event. \yp{} in the macrophage, \ft{} in the macrophage and the
lytic phage in its bacterium are all of that kind, and none is a constant
drip.

That the standard model needs intracellular structure is not a new
observation, and the modelling programme built on it is reviewed by Ciupe and
Conway \cite{ciupe2024incorporating}. Two routes are established. The first
inserts a delay between infection and production, either as a distributed
kernel or as one or more explicit eclipse compartments, an approach going
back to Herz et al.\ \cite{herz1996viral} and now standard
\cite{mittler1998influence,dixit2004complex,kakizoe2015method}. The second,
multiscale, route couples an age-structured population equation to an
explicit intracellular pool: in the hepatitis~C models of Guedj et al.\ and
Rong and Perelson \cite{guedj2013modeling,rong2013mathematical} an
intracellular viral RNA level $R(t,a)$ is tracked alongside infected cells
$I(t,a)$ of infection age $a$, and free virus is fed by
$\rho\int_0^\infty R(t,a)I(t,a)\,\dd a$. Both routes have paid off. The
review's two headline returns are sharper parameter estimates --- neglecting
the pre-productive period biases $R_0$
\cite{lloyd2001dependence} --- and the ability to
locate a drug's mode of action at a named step of the viral life cycle.

What neither route supplies is a \emph{stochastic} intracellular mechanism.
The delay literature assumes a distributional family for the age structure
and fits it; the multiscale literature derives the age structure, but from a
deterministic pool, so it yields a release schedule and nothing more. Yet
for a lytic pathogen the questions that matter early in infection are
distributional. Whether a founder cell establishes an infection at all
depends on how many particles it releases, not on how many it releases on
average, and a cell whose internal population dies out before rupture
releases none.

This chapter closes that gap for one mechanism. \ChCorecap{} and \ChDist{}
developed a birth--death--catastrophe process for the intracellular
replicative units --- birth at rate $\lambda$ per unit, death at rate $\mu$
per unit, and catastrophe at rate $\delta$ per unit, at which the cell
ruptures and releases its entire contents --- and obtained its survival,
release and burst-size laws in closed form. Both stop at the cell boundary.
Here that theory is carried across the boundary into a population.

\subsection*{Contributions}

\begin{enumerate}[leftmargin=1.9em]
\item \textbf{A renewal model forced by the intracellular process}
(\cref{p:sec:renewal}). The population equations need exactly two functions
of cell age: the survival kernel $S$ and the release kernel $g$. Both are
already in closed form, so the population model has no free age structure to
fit. The equations are the first moments of the underlying stochastic
process exactly, not an approximation to them
(\cref{p:prop:first-moment}).

\item \textbf{The standard model is a projection}
(\cref{p:sec:projection}). For each epidemic growth rate $r$ there is a
unique constant pair $\bigl(p_{\rm eff}(r),d_{\Icell,{\rm eff}}(r)\bigr)$
reproducing the renewal system's exponential solution exactly
(\cref{p:thm:projection}). A fitted release rate is therefore a property of
the infection's phase, not of the pathogen: over the whole map it moves by
the factor $\mean{X^2}_{\rm QS}$, a pure number fixed by the intracellular
rates, which is $231$ at $(\lambda,\mu,\delta)=(1,0,0.1)$. No constant pair
serves a transient, and no fit of $(p,d_{\Icell})$ recovers the three
intracellular rates (\cref{p:sec:identifiability}).

\item \textbf{A closed-form answer to whether bursting pays}
(\cref{p:sec:flooding}). This is the result the stochastic treatment exists
to deliver, since it turns on the offspring \emph{distribution} rather than
its mean. Write $L=a(1-b)$ for the flooding parameter
(\cref{p:def:L}). At fixed total output per cell, bursting leaves $R_0$
exactly unchanged (\cref{p:thm:R0inv}); lowers the extinction probability at
establishment if and only if $L>1$, equivalently
\begin{equation*}
\frac{1}{\delta}<\frac{1}{\lambda}+\frac{1}{\mu},
\end{equation*}
by the exact amount $(L-1)(1/m-1)$ (\cref{p:thm:flood}); and slows
deterministic invasion in every regime examined. Threshold, establishment
risk and invasion speed are three questions with three different answers.

\item \textbf{Commensurability} (\cref{p:sec:reach}). Because only the pair
$(S,g)$ enters, budding, bursting, recurrent partial release, multiply
infected cells and the multiscale hepatitis~C model of
\cite{guedj2013modeling,rong2013mathematical} are all instances of one
skeleton, differing only in the kernels they supply
(\cref{p:tab:sg}).
\end{enumerate}

\Cref{p:tab:formulae} collects every population-level formula of the
chapter in one place, for use as a reference card.
Throughout, $\lambda$, $\mu$ and $\delta$ are the intracellular birth, death
and catastrophe rates per replicative unit; the roots $a,b$, the derived
constants $A,B,\theta,\kappa$ and the single-cell functions
$I,D,\Ifix,J,K,V$ are those of \ChCore{}, collected in
\cref{p:app:quoted}. \Cref{p:tab:notation} fixes the population symbols. Two
naming decisions are made once and kept: the burst-time density written
$\varphi$ in \ChDist{} is written $\delta J$ here, because $\varphi$ is
reserved for the partial-release fraction of \cref{p:sec:partial}; and the
survival kernel $S$ is the non-fixation probability $\Ifix$, written $S$
where it acts as a kernel and $\Ifix$ where the closed form is meant.

\begin{table}[htbp]
\centering
\renewcommand{\arraystretch}{1.12}
\small
\begin{tabular}{@{}lp{0.42\textwidth}p{0.30\textwidth}@{}}
\toprule
Symbol & Meaning & Notes \\
\midrule
$\alpha$ & cell age & argument of every kernel \\
$T$, $\gamma$, $c$ & target cells, infection rate, virion clearance & $T$ held fixed \\
$\Icell,\Vfree$ & infected cells, free particles & population counts \\
$i(t)$ & incidence, $\gamma T\Vfree$ & \\
$d_{\Icell}$ & infected-cell removal rate & phenomenological \\
$p$, $p_{\rm eff}(r)$ & constant, effective release rate & \\
$S(\alpha)$, $g(\alpha)$ & survival kernel $\Ifix(\alpha)$, release kernel $\delta K(\alpha)$ & \\
$\mathcal A(\alpha)$ & generation kernel & \cref{p:eq:genkernel} \\
$r$ & epidemic growth rate & \\
$q$, $m$ & infection success $\gamma T/(\gamma T+c)$, branching mean $qV_\infty$ & \\
$L$ & flooding parameter $a(1-b)$ & \cref{p:eq:Ldef} \\
$\varphi$ & partial-release fraction & \cref{p:sec:partial} \\
\bottomrule
\end{tabular}
\caption{Population notation. Single-cell symbols are those of
\cref{p:app:quoted}.}
\label{p:tab:notation}
\end{table}

\subsection{Standing assumptions}
\label{p:sec:assumptions}

These hold throughout unless stated otherwise.

\begin{description}[leftmargin=2.2em,style=sameline]
\item[(A1)] Target cells $T$ are held constant: no depletion, no
saturation, no immune response.
\item[(A2)] Virions are not consumed on infection --- the \emph{continuum}
convention. This is why the deterministic threshold is
$R_0=\gamma TV_\infty/c$ while the discrete branching mean of
\cref{p:sec:flooding} is $m=qV_\infty$ with $q=\gamma T/(\gamma T+c)$; the
two agree when $\gamma T\ll c$.
\item[(A3)] The $\Icell_0$ cells present at $t=0$ all have age exactly
zero.
\item[(A4)] Infected cells are independent, and the population layer
assumes single infection: kernels for founding multiplicity $k$ are
available (\cref{p:app:quoted}) but are not used.
\item[(A5)] There is no spatial structure.
\end{description}

\FloatBarrier

\section{The standard model and the cell it describes exactly}
\label{p:sec:comparator}

\subsection{The standard model}
\label{p:sec:classical}

In the target-cell-constant Basic Model of Viral Replication (BMVR),
infected cells $\Icell$ and free particles $\Vfree$ obey
\begin{equation}
\ddt{\Icell} = \gamma T\,\Vfree - d_{\Icell}\,\Icell,
\qquad
\ddt{\Vfree} = p\,\Icell - c\,\Vfree,
\label{p:eq:bmvr-classical}
\end{equation}
with target cells $T$ fixed, infection rate $\gamma$ and virion clearance
$c$. The basic reproduction number is
\begin{equation}
R_0=\frac{\gamma T\,p}{c\,d_{\Icell}}.
\label{p:eq:R0classical}
\end{equation}
Every cell, whatever its age or contents, releases at rate $p$ and dies at
rate $d_{\Icell}$.

\subsection{The budding cell is described exactly}
\label{p:sec:budding}

This is not an arbitrary caricature. It is one kind of cell described
exactly, and since that cell is the comparator for everything below it is
worth setting out in full. A budding infected cell experiences two
\emph{independent} event streams: release of one particle, as a Poisson
process of rate $p$, and death, as a Poisson process of rate $d_{\Icell}$.
Its survival and cumulative release are
\begin{equation}
S_{\rm bud}(\alpha)=\ee^{-d_{\Icell}\alpha},\qquad
V_{\rm bud}(\alpha)=\frac{p}{d_{\Icell}}\lrb{1-\ee^{-d_{\Icell}\alpha}},
\label{p:eq:bud-cell}
\end{equation}
so the mean productive lifetime is $1/d_{\Icell}$ and the lifetime release
$p/d_{\Icell}$. The total number released is geometric on
$\{0,1,2,\ldots\}$,
\begin{equation}
\pr{\Kb_{\rm bud}=k}=\frac{d_{\Icell}}{d_{\Icell}+p}
\Bigl(\frac{p}{d_{\Icell}+p}\Bigr)^{\!k},
\label{p:eq:bud-geom}
\end{equation}
the successes before the terminating death in a sequence of events each of
which is a release with probability $p/(d_{\Icell}+p)$. All of the cell sits
in the pair $(p,d_{\Icell})$. That is why \cref{p:eq:bmvr-classical} is a
faithful one-cell limit for a budding pathogen, and it is this cell that
\cref{p:sec:flooding} matches against a bursting one.

\subsection{Bursting, and why no constant repairs it}
\label{p:sec:obstruction}

In bursting, release and death are the same event: whatever has accumulated
leaves at once and the cell's productive life ends in that instant. Take the
cleanest consequence. Let $N_0$ cells be infected at the same moment, each
running an independent copy of the single-cell process. The number still
productively infected at time $t$ is $N_0\Ifix(t)$ exactly in expectation,
and the expected cumulative release is $N_0V(t)$, so the release
\emph{rate} is $N_0\delta K(t)$. It depends on the age of the cohort and it
is not constant. A cohort releases on a schedule, not at a rate
(\cref{p:fig:constant-fails}).

\begin{figure}[tbp]
\centering
\includegraphics[width=0.94\textwidth]{N4b_1_constant_release_fails.pdf}
\caption{A synchronous cohort releases on a schedule, not at a rate. Both
schedules have lifetime yield $V_\infty=11$ and mean productive lifetime
$\mean{T_{\rm prod}}=2.398$, with $d_{\Icell}=1/\mean{T_{\rm prod}}$,
$p=V_\infty d_{\Icell}$, $(\lambda,\mu,\delta)=(1,0,0.1)$ and $c=1$: as
fairly matched as two constants allow. (a) At matched yield the release
kernel $g(\alpha)=\delta K(\alpha)$ still lags the budding cell's
$p\ee^{-d_{\Icell}\alpha}$, because a bursting cell must first accumulate
load. (b) Linear clearance preserves that lag, carrying it into the
free-particle response.}
\label{p:fig:constant-fails}
\end{figure}

The natural repair is to keep \cref{p:eq:bmvr-classical} and rebuild $p$
from the intracellular theory --- to set $p$ equal to the quasi-stationary
mean load $\mean{X}_{\rm QS}=a/(a-1)$, on the reasoning that a long-lived
cell hovers at that level. The obstruction is dimensional. Here $p$ is a
number of particles per unit time and $\mean{X}_{\rm QS}$ is a pure count,
so the substitution introduces a spurious factor of time; when $\mu=0$ one
has $\mean{X}_{\rm QS}=V_\infty$ exactly, and the proposal replaces a cell's
output \emph{rate} by its total lifetime \emph{output}. Averaging the burst
size over the burst-time density fails the same way, returning
$\delta\int_0^\infty K=V_\infty$, a count again. Weighting the
load-dependent flux by the burst-time density has the right dimensions but
the wrong measure: that density says when a \emph{single} cell bursts,
whereas a population release rate needs an average over the ages of the
cells \emph{currently infected}, and in an asynchronous infection those are
not the same.

What survives is two facts rather than a substitute constant. The expected
release flux of a cell of age $\alpha$ is $\delta K(\alpha)$; and at any
host time the infected population is a mixture of ages, one cohort per past
moment of infection. The population model must therefore track cell age.

\section{The intracellular process and its two kernels}
\label{p:sec:kernels}

\subsection{The process}

Inside a cell, replicative units follow a linear birth--death process with
an added catastrophe. Each unit gives birth at rate $\lambda$, dies at rate
$\mu$, and at rate $\delta$ per unit triggers rupture, at which the cell
releases its entire load and enters an absorbing state. Writing $X_\alpha$
for the load at cell age $\alpha$, the process is absorbed either by
rupture --- a burst --- or by internal extinction, in which the last unit
dies and the cell never bursts at all. \ChCorecap{} solves it. The roots
\begin{equation}
a=\eta+\sqrt{\eta^2-\tfrac{\mu}{\lambda}},\qquad
b=\eta-\sqrt{\eta^2-\tfrac{\mu}{\lambda}},\qquad
\eta=\frac{\lambda+\mu+\delta}{2\lambda},
\label{p:eq:roots}
\end{equation}
satisfy $b<1<a$ and carry everything: $b$ is the probability that a cell
clears internally without ever bursting, and $1/a$ is the ratio of the
geometric burst-size law. \Cref{p:app:quoted} collects the closed forms used
below; the two that matter most are the lifetime yield and the mean
productive lifetime,
\begin{equation}
V_\infty=\frac{a(1-b)}{a-1},
\qquad
\mean{T_{\rm prod}}=\frac{1}{\lambda}\log\frac{a}{a-1}.
\label{p:eq:Vinf-Tprod}
\end{equation}

\subsection{The two kernels}

Age-structured population dynamics needs exactly two single-cell
ingredients. The \emph{survival kernel} is the probability that a cell
infected $\alpha$ time units ago is still productively infected,
\begin{equation}
S(\alpha):=\Ifix(\alpha)=\frac{(a-b)^2\,\ee^{\theta\alpha}}
{\lrb{B+A\ee^{\theta\alpha}}\lrb{a\ee^{\theta\alpha}-b}},
\label{p:eq:Skernel}
\end{equation}
and the \emph{release kernel} is the expected release flux at cell age
$\alpha$,
\begin{equation}
g(\alpha):=\delta K(\alpha)=V'(\alpha).
\label{p:eq:gkernel}
\end{equation}
The reasoning behind \cref{p:eq:gkernel} is the one the failed proposal of
\cref{p:sec:obstruction} stated and then did not use: a cell bursts at rate
$\delta X_\alpha$ and releases $X_\alpha$ particles when it does, so the
expected number released in $[\alpha,\alpha+\dd\alpha]$ is
$\delta\mean{X_\alpha^2}\dd\alpha=\delta K(\alpha)\dd\alpha$. With $K$ from
\cref{p:eq:K} both kernels are explicit functions of age
(\cref{p:fig:kernels}). This is what it will mean for the population model
to be \emph{forced by} the intracellular process: the kernels are computed,
not fitted.

\begin{figure}[tbp]
\centering
\includegraphics[width=0.98\textwidth]{N4b_7_kernels_three_regimes.pdf}
\caption{Survival $\Ifix$ (a)--(c) and release $g=\delta K$ (d)--(f) across
the flooding boundary of \cref{p:tab:trio}. Release rises then falls --- a
cell must accumulate before it bursts --- while survival decays; shaded
areas check $\int_0^\infty\Ifix=\mean{T_{\rm prod}}$ and
$\int_0^\infty g=V_\infty$. Parameters, left to right:
$(\lambda,\mu,\delta)=(1,0,0.1)$, $(1,0.5,1/3)$, $(1,0.9,0.1)$, with
$L=1.10$, $1.00$, $0.42$ and $V_\infty=11.00$, $2.00$, $1.32$. Independent
$g$-scales keep the timing shift visible.}
\label{p:fig:kernels}
\end{figure}

\subsection{Three regimes, used throughout}
\label{p:sec:trio}

Three intracellular parameter triples recur below, chosen to straddle the
flooding boundary $L=1$ of \cref{p:sec:flooding}: one on each side, one
exactly on it. They are named here and used by name thereafter.

\begin{table}[htbp]
\centering
\renewcommand{\arraystretch}{1.12}
\small
\begin{tabular}{@{}l
    S[table-format=1.3]S[table-format=1.3]S[table-format=1.2]
    S[table-format=2.2]S[table-format=1.3]S[table-format=1.3]
    S[table-format=1.2]@{}}
\toprule
{$(\lambda,\mu,\delta)$} & {$a$} & {$b$} & {$L$} & {$V_\infty$}
 & {$\mean{T_{\rm prod}}$} & {$d_{\Icell,{\rm eff}}(0)$} & {$p_{\rm eff}(0)$} \\
\midrule
$(1,0,0.1)$    & 1.100 & 0     & 1.10 & 11.00 & 2.398 & 0.417 & 4.59 \\
$(1,0.5,1/3)$  & 1.500 & {$1/3$} & 1.00 & 2.00  & 1.099 & 0.910 & 1.82 \\
$(1,0.9,0.1)$  & 1.316 & 0.684 & 0.42 & 1.32  & 1.426 & 0.701 & 0.92 \\
\bottomrule
\end{tabular}
\caption{The three regimes used throughout. The first has $L>1$, the second
sits on the boundary $L=1$, the third has $L<1$. The last two columns are
the $r=0$ endpoints of the effective-parameter map of
\cref{p:sec:projection}.}
\label{p:tab:trio}
\end{table}

\FloatBarrier

\section{Burst-aware renewal dynamics}
\label{p:sec:renewal}

\subsection{The renewal system}

Hold target cells $T$ fixed, let infection occur at rate $\gamma$ per
virion, clear virions at rate $c$, and define the incidence
$i(t)=\gamma T\Vfree(t)$. Each infection event at time $t-\alpha$ starts a
cell of age $\alpha$, still productive with probability $S(\alpha)$ and
releasing at expected rate $g(\alpha)$. Summing over infection times and
adding the $\Icell_0$ cells already infected at $t=0$ gives the system,
whose bookkeeping is drawn in \cref{p:fig:renewal-schematic}.

\begin{definition}[The burst-aware renewal system]
\label{p:def:renewal-system}
\begin{equation}
\Icell(t)=\Icell_0\,\Ifix(t)+\int_0^t i(t-\alpha)\,\Ifix(\alpha)\,\dd\alpha
=\Icell_0\,\Ifix+(i*\Ifix)(t),
\label{p:eq:Icell}
\end{equation}
\begin{equation}
\ddt{\Vfree}(t)
=\Icell_0\,g(t)+\int_0^t i(t-\alpha)\,g(\alpha)\,\dd\alpha-c\,\Vfree(t)
=\Icell_0\,g+(i*g)(t)-c\,\Vfree(t).
\label{p:eq:Vfree}
\end{equation}
\end{definition}

\begin{figure}[t]
\centering
\includegraphics[width=0.98\textwidth]{F4b_1_renewal_schematic.pdf}
\caption{Bookkeeping of \cref{p:def:renewal-system}: infection events start
cohorts; the kernels $\Ifix$ and $g=\delta K$ weight them by age; the
feedback $i=\gamma T\Vfree$ closes the loop.}
\label{p:fig:renewal-schematic}
\end{figure}

\Cref{p:tab:correspondence} sets the two systems side by side. Both
constants are replaced, not only the release: memoryless removal
$d_{\Icell}\Icell$ becomes the age structure implicit in \cref{p:eq:Icell},
and constant release $p\Icell$ becomes the convolution with $g$. A
replacement touching only the release term --- the first thought of
\cref{p:sec:obstruction} --- would leave the removal side memoryless, and
still wrong. The removal side is not merely a formal concern: analysing
stage-resolved HIV decay under drugs acting at different points in the life
cycle, Gilmore et al.\ \cite{gilmore2013explaining} found infected-cell
half-lives close to those of uninfected cells before integration and much
shorter after it, which is an age-dependent $d_{\Icell}$ measured directly.

\begin{table}[htbp]
\centering
\renewcommand{\arraystretch}{1.25}
\small
\begin{tabular}{@{}p{0.46\textwidth}p{0.48\textwidth}@{}}
\toprule
Standard BMVR \cref{p:eq:bmvr-classical}
& Renewal system \crefrange{p:eq:Icell}{p:eq:Vfree} \\
\midrule
$\Icell'=i-d_{\Icell}\Icell$
& $\Icell=\Icell_0\Ifix+i*\Ifix$ \\
constant removal rate $d_{\Icell}$
& age-structured survival via $\Ifix$ \\[0.5em]
$\Vfree'=p\Icell-c\Vfree$
& $\Vfree'=\Icell_0 g+i*g-c\Vfree$ \\
constant release rate $p$
& age-structured release via $g=\delta K$ \\
\bottomrule
\end{tabular}
\caption{What the renewal system replaces. Both constants go, not only the
release rate.}
\label{p:tab:correspondence}
\end{table}

Differentiating \cref{p:eq:Icell} makes the structure visible:
\begin{equation*}
\ddt{\Icell}
= i(t)+\Icell_0\,\Ifix'(t)+\int_0^t i(t-\alpha)\,\Ifix'(\alpha)\,\dd\alpha,
\end{equation*}
in which incidence appears as a source exactly as in the standard model,
while the remaining terms encode age-dependent loss of productivity. There
is in general no constant $d_{\Icell}$ reducing this to
$\Icell'=i-d_{\Icell}\Icell$ for every history of $i$; that holds only when
everything grows exponentially, and the exception is what
\cref{p:sec:projection} exploits.

\subsection{The equations are exact in the mean}

Equations \cref{p:eq:Icell,p:eq:Vfree} are not an approximation to the
underlying stochastic process. They are its first moments exactly, and the reason is
that every rate in sight is linear.

\begin{proposition}[First-moment exactness]
\label{p:prop:first-moment}
Consider the stochastic process in which each infection event starts an
independent birth--death--catastrophe cell, each released particle
independently infects at rate $\gamma T$ or clears at rate $c$, and the
$\Icell_0$ cells present at $t=0$ have age zero. Write $\Icell(t)$ and
$\Vfree(t)$ for the expected numbers of productive cells and free particles.
Then $\Icell$ and $\Vfree$ satisfy \cref{p:eq:Icell,p:eq:Vfree} exactly.
\end{proposition}

\begin{proof}
Let $N$ be the point process of infection events. Since each free particle
infects at rate $\gamma T$ and particles are independent, $N$ has
compensator $\gamma T\,\Vfree^{\omega}(t)\,\dd t$, where $\Vfree^\omega$ is
the realised particle count; taking expectations, its mean intensity is
$i(t)=\gamma T\Vfree(t)$. A cell founded at time $t-\alpha$ evolves
independently of $N$ and of every other cell, so it is still productive at
$t$ with probability $\Ifix(\alpha)$ and its expected release flux at $t$ is
$g(\alpha)$, both by definition of the kernels. Conditioning on $N$ and
applying Campbell's formula --- legitimate because $\Ifix$ and $g$ are
bounded and non-negative, so Fubini applies --- gives
$\Icell(t)=\Icell_0\Ifix(t)+(i*\Ifix)(t)$ and an expected release flux of
$\Icell_0g(t)+(i*g)(t)$. Clearance removes particles at rate $c$ per
particle, contributing $-c\Vfree(t)$ to the mean. Collecting terms gives
\cref{p:eq:Vfree}.
\end{proof}

\noindent
The step that would fail for a nonlinear incidence is the identification of
the mean intensity with $\gamma T\Vfree$: for $f(\Vfree,T)$ nonlinear,
$\mean{f(\Vfree^\omega)}\neq f(\mean{\Vfree^\omega})$ and the moment
hierarchy does not close. Assumption~(A1) is what keeps the incidence
linear. This is one reason to prefer mass action where the data permit it,
and the review of \cite{ciupe2024incorporating} reports that they generally
do, the saturating alternatives --- Holling type~2, Beddington--DeAngelis,
Crowley--Martin --- being well explored theoretically but not validated
against data.

\subsection{One skeleton, many mechanisms}
\label{p:sec:generality}

Nothing in \cref{p:eq:Icell,p:eq:Vfree} used the birth--death--catastrophe
process beyond the two kernels it supplied. With incidence
$i(t)=\gamma T\Vfree(t)$, or a nonlinear $f(\Vfree,T)$, \emph{any} pair
$(S,g)$ gives
\begin{equation}
\Icell(t)=\Icell_0\,S(t)+(i*S)(t),
\qquad
\ddt{\Vfree}=\Icell_0\,g(t)+(i*g)(t)-c\,\Vfree(t),
\label{p:eq:skeleton}
\end{equation}
with effective parameters $p_{\rm eff}(r)=\Lap{g}(r)/\Lap{S}(r)$ and, for
linear incidence, $R_0=(\gamma T/c)\int_0^\infty g$. This is the chapter's
most general statement, and it is what makes release mechanisms
commensurable: they differ only in $(S,g)$ (\cref{p:tab:sg}). The first row
is a structural check on the whole construction: exponential kernels reduce
\cref{p:eq:skeleton} exactly to standard BMVR
(\cref{p:sec:verification}).

\begin{table}[htbp]
\centering
\renewcommand{\arraystretch}{1.1}
\small
\begin{tabular}{@{}p{0.26\textwidth}p{0.19\textwidth}p{0.41\textwidth}@{}}
\toprule
Model & $S(\alpha)$ & $g(\alpha)$ \\
\midrule
Standard BMVR \cite{perelson1996hiv} & $\ee^{-d_{\Icell}\alpha}$ & $p\,\ee^{-d_{\Icell}\alpha}$ \\
Eclipse $+$ standard \cite{ciupe2024incorporating} & $\ee^{-d_{\Icell}\alpha}$ & $p\ee^{-d_{\Icell}\alpha}\ind\{\alpha>\tau_e\}$ \\
HCV multiscale \cite{guedj2013modeling,rong2013mathematical} & $\ee^{-\delta\alpha}$ & $\rho R(\alpha)\ee^{-\delta\alpha}$, with $R$ from an intracellular ODE \\
Absorbing BDC (this chapter) & $\Ifix(\alpha)$ & $\delta K(\alpha)$ \\
MOI-$k$ BDC (\ChDist) & $I(\alpha)^k-D(\alpha)^k$ & $\delta K_k(\alpha)$ \\
Reset catastrophe (\cref{p:sec:reset}) & $\ee^{-d_{\Icell}\alpha}$ & $\delta\ee^{-d_{\Icell}\alpha}\mean{X_\alpha^2}$ \\
Two-type / GATE (\ChTwoType) & two-type $\Ifix$ & two-type second-moment flux \\
\bottomrule
\end{tabular}
\caption{One skeleton, many mechanisms. Each row is a different single-cell
story; all induce \cref{p:eq:skeleton}, and differ only in the pair of
kernels they supply.}
\label{p:tab:sg}
\end{table}

\subsection{Relation to existing multiscale models}
\label{p:sec:relation}

The renewal framework itself is standard: it is the McKendrick--von Foerster
machinery of mathematical demography
\cite{mckendrick1926applications,vonfoerster1959some,metz1986dynamics}. So
are age-structured within-host models. Nelson et al.\
\cite{nelson2004agestructured} write down a within-host system with
age-dependent production and death rates that is structurally
\cref{p:eq:skeleton}, and take the age dependence as given. The multiscale
hepatitis~C models \cite{guedj2013modeling,rong2013mathematical} go further
and derive it, coupling infected cells of age $a$ to an intracellular viral
RNA level $R(t,a)$ obeying a transport equation with production, degradation
and export; free virus is then fed by
$\rho\int_0^\infty R(t,a)I(t,a)\,\dd a$, which is
\cref{p:eq:skeleton} with $g(\alpha)=\rho R(\alpha)\ee^{-\delta\alpha}$.
That model is the third row of \cref{p:tab:sg} and a direct ancestor of what
is done here.

The difference is what the intracellular layer is. In
\cite{guedj2013modeling,rong2013mathematical} it is a deterministic pool, so
it supplies a mean release schedule and nothing else. Here it is a stochastic
process solved in closed form, so it supplies the schedule \emph{and} the
distributions built on it: the burst-size law, the probability $b$ that a
cell clears internally and releases nothing, the offspring generating
function, and hence extinction probabilities. Everything in
\cref{p:sec:projection} could in principle be asked of a deterministic
intracellular model; nothing in \cref{p:sec:flooding} could, because
establishment is a question about the offspring distribution rather than its
mean. That is the sense in which the stochastic mechanism earns its place.

\subsection{Numerical verification}
\label{p:sec:verification}

The renewal system was checked end to end at four levels: kernel integrity,
the reduction at $\gamma T=0$, the structural limits available in closed
form, and stochastic ground truth. All 62 checks pass. Exponential kernels
reproduce standard BMVR to relative error below $5\times10^{-5}$; the late
growth rate equals the characteristic root of \cref{p:sec:projection} to
about $10^{-7}$; the $R_0$ threshold is sharp; and an independent Gillespie
ensemble reproduces the renewal means for $\Icell$ and $\Vfree$ to within
$1$--$2\%$ (\cref{p:fig:gillespie}). The suite runs over the three regimes
of \cref{p:tab:trio} together with $(1,0.2,0.05)$ and $(1,0,0.5)$, chosen to
spread the kernel shapes; it is the script
\texttt{verify\_result\_20\_1.py} distributed with the chapter source, and
tests~M and~N of it are the two numerical sweeps cited in
\cref{p:sec:span,p:sec:tradeoff}.
% AUTHOR-ACTION: replace with a repository URL or archived DOI when one
% exists.

\begin{figure}[tbp]
\centering
\includegraphics[width=0.98\textwidth]{H_gillespie.pdf}
\caption{Stochastic ground truth. Ensemble means of an independent Gillespie
simulation --- birth--death--catastrophe cells with impulse releases at
burst times, exponential clearance, no reinfection --- against
\crefrange{p:eq:Icell}{p:eq:Vfree}. Left column $\mu=0$, right column
$\mu>0$; (a)--(b) infected cells, (c)--(d) free particles. The curves are
indistinguishable at plotting resolution and the relative $L^2$ discrepancy
of $1$--$2\%$ is sampling error: the renewal equations are not an
approximation to the simulation, and this is the check that says so.}
\label{p:fig:gillespie}
\end{figure}

\FloatBarrier

\section{The standard model is a projection}
\label{p:sec:projection}

The renewal system does not merely disagree with standard BMVR in
transients. It contains standard BMVR as an exact special case, and the
relationship is precise enough to be turned into predictions.

\subsection{Exponential-phase equivalence}

Suppose the incidence grows exponentially, $i(t)=i_0\ee^{rt}$, at times late
enough that the lower limits of the convolutions are irrelevant. Then
\begin{equation}
\Icell(t)=i_0\ee^{rt}\Lap{\Ifix}(r),
\qquad
\Vfree'(t)=i_0\ee^{rt}\,\delta\Lap{K}(r)-c\,\Vfree(t),
\label{p:eq:expphase}
\end{equation}
where $\Lap{f}(r)=\int_0^\infty\ee^{-r\alpha}f(\alpha)\dd\alpha$. A standard
model driven by the same incidence has $\Icell(t)=i_0\ee^{rt}/(r+d_{\Icell})$
and the eigenvalue relation $(r+d_{\Icell})(r+c)=\gamma Tp$; matching the two
exponential solutions fixes
\begin{equation}
p_{\rm eff}(r)=\frac{\delta\,\Lap{K}(r)}{\Lap{\Ifix}(r)},
\qquad
d_{\Icell,{\rm eff}}(r)=\frac{1}{\Lap{\Ifix}(r)}-r.
\label{p:eq:peff}
\end{equation}

\begin{theorem}[Exponential-phase equivalence]
\label{p:thm:projection}
Fix a growth rate $r>-\theta$. The constant pair
$\bigl(p_{\rm eff}(r),d_{\Icell,{\rm eff}}(r)\bigr)$ of \cref{p:eq:peff} is
the unique one for which standard BMVR reproduces the renewal system's
$\ee^{rt}$ solution exactly, in both $\Icell$ and $\Vfree$; and substituting
it into the standard eigenvalue relation returns
\begin{equation}
r+c=\gamma T\,\delta\Lap{K}(r),
\label{p:eq:char}
\end{equation}
precisely the characteristic equation of the renewal system. What no
constant pair reproduces is a transient, which samples a continuum of growth
rates at once.
\end{theorem}

\begin{proof}
Uniqueness and existence are \cref{p:eq:expphase} read as two equations in
two unknowns: matching $\Icell$ forces $1/(r+d_{\Icell})=\Lap{\Ifix}(r)$,
which is $d_{\Icell,{\rm eff}}$, and matching the release term then forces
$p\Lap{\Ifix}(r)=\delta\Lap{K}(r)$, which is $p_{\rm eff}$. Both transforms
are finite and strictly positive for $r>-\theta$, so the pair is well
defined. Substituting into $(r+d_{\Icell})(r+c)=\gamma Tp$ gives
$(r+c)/\Lap{\Ifix}(r)=\gamma T\delta\Lap{K}(r)/\Lap{\Ifix}(r)$, that is
\cref{p:eq:char}. That no single pair serves a transient follows from
\cref{p:prop:peff-monotone}: $p_{\rm eff}$ is not constant in $r$.
\end{proof}

\noindent
So the standard model is not a rival to the renewal system but a projection
of it, exact on one exponential phase at a time. Everything else in this
section reads \cref{p:eq:peff} as a map.

\subsection{The ends of the map}

Using $\int_0^\infty\delta K=V_\infty$ and
$\int_0^\infty\Ifix=\mean{T_{\rm prod}}$,
\begin{equation}
r=0:\quad
p_{\rm eff}(0)=\frac{V_\infty}{\mean{T_{\rm prod}}}
=\frac{\lambda\,V_\infty}{\log\bigl(a/(a-1)\bigr)},
\quad
d_{\Icell,{\rm eff}}(0)=\frac{\lambda}{\log\bigl(a/(a-1)\bigr)},
\label{p:eq:peff0}
\end{equation}
lifetime output divided by lifetime --- a consistency check against naive
moment-matching, and the matching used in the overlays below. At the other
end,
\begin{equation}
r\to\infty:\qquad
p_{\rm eff}(r)\to\delta,
\qquad
d_{\Icell,{\rm eff}}(r)\to\mu+\delta,
\label{p:eq:peffinf}
\end{equation}
the \emph{young-cell} limit: in a fast-growing epidemic almost every
infected cell was infected moments ago, holds one particle, bursts at rate
$\delta$ and releases one.

Calling these the two ends understates the map, because $r=0$ is not an end.
Both transforms converge for every $r>-\theta$, and negative $r$ is not a
fiction --- it is a decaying infection, the case of a treated patient. The
genuine left-hand endpoint mirrors the young-cell limit.

\begin{proposition}[The old-cell limit]
\label{p:prop:oldcell}
The release flux per surviving cell tends to a finite limit,
\begin{equation}
\lim_{\alpha\to\infty}\frac{g(\alpha)}{S(\alpha)}
=\frac{a\lrbs{2\lambda(1-b)+\delta}}{a-1}
=\frac{\delta\,a(a+1)}{(a-1)^2}
=\delta\,\mean{X^2}_{\rm QS},
\label{p:eq:oldcell}
\end{equation}
the burst rate times the second moment of the quasi-stationary load, and
consequently $p_{\rm eff}(r)\to\delta\mean{X^2}_{\rm QS}$ as
$r\downarrow-\theta$.
\end{proposition}

\begin{proof}
Write $v=\ee^{-\theta\alpha}$. As $\alpha\to\infty$, $v\to0$ and
\cref{p:eq:IDIhat,p:eq:J} give $\Ifix\sim(a-b)^2v/(Aa)$ and $I\to b$, so by
\cref{p:eq:K} $g=2\lambda(\kappa-I)J\sim2\lambda(\kappa-b)(a-b)^2v/A^2$. The
ratio of the two is $2\lambda(\kappa-b)a/A$, which is the first equality
after $2\lambda(\kappa-b)=2\lambda(1-b)+\delta$. For the second, the
identity $AB=\delta/\lambda$ reads $\lambda(1-b)=\delta/(a-1)$, whence
$a[2\delta/(a-1)+\delta]/(a-1)=\delta a(a+1)/(a-1)^2$. That last expression
is $\delta\mean{X^2}_{\rm QS}$ because the quasi-stationary load is
geometric on $\{1,2,\ldots\}$ with ratio $1/a$ (\cref{p:eq:burstlaw}), whose
second moment is $a(a+1)/(a-1)^2$. Finally, both $\Lap{\Ifix}$ and
$\delta\Lap K$ diverge like $1/(r+\theta)$ as $r\downarrow-\theta$, with
coefficients in the ratio just computed, so their quotient tends to it.
\end{proof}

\noindent
The limit \cref{p:eq:oldcell} also repairs the failed proposal of
\cref{p:sec:obstruction}, which offered the quasi-stationary \emph{mean} as
a release rate and was wrong twice over: a count where a rate was needed,
and the first moment where the second is called for. Both corrections appear
together --- $\delta$ supplies the missing inverse time, and the
load-proportional burst hazard weights by $X$ a second time.

\subsection{Monotonicity, and the size of the effect}
\label{p:sec:span}

\begin{proposition}[Monotonicity of the effective release rate]
\label{p:prop:peff-monotone}
Write $\rho(\alpha)=g(\alpha)/S(\alpha)$ for the expected release flux per
still-productive cell, so that $\rho(0)=\delta$ and
$\rho(\infty)=\delta\mean{X^2}_{\rm QS}$. If $\rho$ is non-decreasing then
$p_{\rm eff}$ is non-increasing on $(-\theta,\infty)$, strictly so if $\rho$
is strictly increasing.
\end{proposition}

\begin{proof}
For $r>-\theta$ let $\pi_{f,r}$ be the probability density
$\alpha\mapsto\ee^{-r\alpha}f(\alpha)/\Lap f(r)$. Differentiating under the
integral, $\dd\log\Lap f(r)/\dd r=-\mean{\alpha}_{\pi_{f,r}}$, so
\begin{equation}
\frac{\dd}{\dd r}\log p_{\rm eff}(r)
=\mean{\alpha}_{\pi_{S,r}}-\mean{\alpha}_{\pi_{g,r}}.
\label{p:eq:peff-deriv}
\end{equation}
The two densities have ratio
$\pi_{g,r}(\alpha)/\pi_{S,r}(\alpha)\propto\rho(\alpha)$, independently of
$r$. If $\rho$ is non-decreasing, $\pi_{g,r}$ dominates $\pi_{S,r}$ in the
monotone likelihood ratio order, hence in the usual stochastic order, hence
in mean; so \cref{p:eq:peff-deriv} is non-positive.
\end{proof}

\noindent
The hypothesis is a statement about one closed-form function of one
variable: substituting $v=\ee^{-\theta\alpha}$ gives
\begin{equation}
\rho=\frac{2\lambda\lrbs{B(\kappa-a)v+A(\kappa-b)}(a-bv)}{(Bv+A)^2},
\label{p:eq:rho-closed}
\end{equation}
which must be shown decreasing in $v\in(0,1]$. That reduces the question to
a finite algebraic one. It is not settled here, but $\rho$ was found
increasing at every one of $4000$ triples sampled over
$(\mu,\delta)\in[10^{-3},20]^2$ at $\lambda=1$, and $p_{\rm eff}$ decreasing
at every one (test~M of the verification suite).

The size of the effect is what matters. At $(\lambda,\mu,\delta)=(1,0,0.1)$,
$p_{\rm eff}(0)\approx4.59$ against $p_{\rm eff}(\infty)=0.1$: the same
cell, two orders of magnitude apart, with nothing separating the two figures
but how fast the epidemic around it is growing. Over the whole map the span
is larger, and exact --- by \cref{p:eq:peffinf,p:eq:oldcell} the ratio of
the endpoints is
\begin{equation}
\frac{\delta\,\mean{X^2}_{\rm QS}}{\delta}=\mean{X^2}_{\rm QS}
=\frac{a(a+1)}{(a-1)^2},
\label{p:eq:peff-span}
\end{equation}
a pure number fixed by the intracellular rates alone. It is $231$ at
$(1,0,0.1)$ and $30.5$ at $(1,0.9,0.1)$: \emph{the full dynamic range of a
fitted release rate is the second moment of the quasi-stationary load.}

\begin{figure}[tbp]
\centering
\includegraphics[width=0.94\textwidth]{peff_dr_curves.pdf}
\caption{The effective release rate (a) and removal rate (b) of
\cref{p:eq:peff} against the epidemic growth rate $r$, for three parameter
sets, over the whole domain $r>-\theta$ on which the map is defined. Open
circles mark the $r=0$ values \cref{p:eq:peff0}; dotted segments mark the
limits each curve runs toward. Note that $r=0$ is an interior point: the
left-hand endpoint is the old-cell limit $\delta\mathbb{E}_{\rm QS}[X^2]$
and the right-hand one the young-cell limit $\delta$. Read the vertical
extent of each curve as the size of the error made by treating $p$ as a
constant of the pathogen rather than of the epidemic's phase.}
\label{p:fig:peff}
\end{figure}

\subsection{Threshold and growth rate as functionals of one kernel}
\label{p:sec:generation-kernel}

One object lies behind both \cref{p:eq:peff} and the comparison of
\cref{p:sec:flooding}. Define the cell-to-cell \emph{generation kernel}
\begin{equation}
\mathcal A(\alpha)
=\gamma T\,\bigl(g*\ee^{-c\,\cdot}\bigr)(\alpha)
=\gamma T\!\int_0^\alpha g(u)\,\ee^{-c(\alpha-u)}\,\dd u,
\label{p:eq:genkernel}
\end{equation}
the rate at which a cell of age $\alpha$ generates new infected cells: it
releases at rate $g(u)$ at age $u$, a particle released then is still free a
time $\alpha-u$ later with probability $\ee^{-c(\alpha-u)}$, and infects at
rate $\gamma T$. Its transform is
$\Lap{\mathcal A}(r)=\gamma T\Lap{g}(r)/(r+c)$, so the characteristic
equation \cref{p:eq:char} is exactly $\Lap{\mathcal A}(r)=1$: the
Euler--Lotka equation of a one-type age-structured population, with the
infected cell as the individual and $\mathcal A$ as its mean reproduction
measure \cite{diekmann1990definition,metz1986dynamics}.

Three things follow. First, the total mass of $\mathcal A$ is the basic
reproduction number,
\begin{equation}
R_0=\int_0^\infty\mathcal A(\alpha)\,\dd\alpha
=\gamma T\int_0^\infty g\;\cdot\;\int_0^\infty\ee^{-cu}\,\dd u
=\frac{\gamma T\,V_\infty}{c},
\label{p:eq:R0gen}
\end{equation}
which depends on $g$ only through its integral. Second, the normalised
kernel $\mathcal A/R_0$ is the \emph{generation-time density}, the law of
the age at which a cell produces a randomly chosen one of its offspring
cells. Third, since $r$ solves $\Lap{\mathcal A}(r)=1$ while
$R_0=\Lap{\mathcal A}(0)$, threshold and growth rate are different
functionals of one kernel --- the standard $R_0$-versus-$r$ distinction
\cite{wallinga2007generation}, here with both sides in closed form. The
first two of these deliver the invariance that makes
\cref{p:sec:flooding} a fair comparison.

\begin{theorem}[$R_0$ is unchanged by bursting]
\label{p:thm:R0inv}
At fixed total per-cell output $V_\infty$, replacing constant budding by
bursting leaves $R_0$ exactly unchanged. Bursting alters the generation-time
distribution --- hence the growth rate $r$, hence the extinction
probability, hence the stochastic early phase --- but not the threshold
itself.
\end{theorem}

\begin{proof}
By \cref{p:eq:R0gen}, $R_0=\int_0^\infty\mathcal A$ depends on the release
kernel only through $\int_0^\infty g=V_\infty$. Replacing bursting by
budding at matched yield replaces $g$ by another kernel with the same
integral, and leaves $R_0$ where it was. The generation-time density
$\mathcal A/R_0$, by contrast, depends on the whole of $g$, and it is
through that density that $r$ is determined by $\Lap{\mathcal A}(r)=1$.
\end{proof}

\noindent
A caution attaches about assumption (A2): \cref{p:eq:R0gen} is the
\emph{continuum} reproduction number, in which one virion contributes mean
$\gamma T/c$ new infections, whereas a discrete branching process absorbing
virions on infection carries the mean $m=qV_\infty$ with
$q=\gamma T/(\gamma T+c)$. The two agree when $\gamma T\ll c$, and it is $m$
that governs \cref{p:sec:flooding}.

\subsection{What the projection loses}
\label{p:sec:overlays}

To see what age structure does to an epidemic, integrate both systems side
by side. In \cref{p:fig:overlay-V,p:fig:growth-phase} the solid curve is the
renewal system and the dashed one standard BMVR with the fairest constant
match, the $r=0$ end of \cref{p:eq:peff}, so that total lifetime output and
mean productive lifetime agree.

\begin{figure}[tbp]
\centering
\includegraphics[width=0.94\textwidth]{overlay_V.pdf}
\caption{Both state variables, both sides of the threshold. Free particles
$\Vfree(t)$ (top row) and infected cells $\Icell(t)$ (bottom) under the
renewal system (solid) and under standard BMVR matched at $r=0$ (dashed),
for a supercritical infection (left, $R_0=1.76$) and a subcritical one
(right, $R_0=0.28$). The matched model gets the order of magnitude and not
the shape; and the discrepancy is not confined to the release equation ---
age structure moves the cell trajectory too, which is why matching $p$ alone
cannot repair it. Parameters: $(\lambda,\mu,\delta)=(1,0,0.1)$ throughout,
with $\gamma T=0.04$, $c=0.25$ on the left and $\gamma T=0.01$, $c=0.4$ on
the right.}
\label{p:fig:overlay-V}
\end{figure}

The reading is consistent across panels. The matched model misses three
things: the early lag, because young cells dominate and release is nearer
the young-cell limit $\delta$ than $p_{\rm eff}(0)$; the timing of the peak;
and the late growth rate, since generation-time structure moves the
operative release rate away from its $r=0$ value as the epidemic grows. The
relative free-particle discrepancy is often order one or larger by the end of
the plotted window in supercritical regimes, and worst when intracellular
death is high, where $V_\infty$ is small while the age structure of release
stays strongly non-exponential.

\begin{figure}[tbp]
\centering
\includegraphics[width=0.94\textwidth]{overlay_growth_phase.pdf}
\caption{Why no single constant pair will do. One renewal trajectory (solid)
against two standard parameter choices: the $r=0$ match (dashed) and the
young-cell limit $p=\delta$, $d_{\Icell}=\mu+\delta$ (dash-dot). (a) The
first six time units, where the $r=0$ match overshoots by more than an order
of magnitude while the young-cell parameters are the closer of the two.
(b) The full window, where the ordering has reversed and the young-cell
model has collapsed. Each choice is correct at one growth rate and the
trajectory passes through many, exactly as \cref{p:eq:peff} predicts. Both
panels logarithmic in $\Vfree$. Parameters:
$(\lambda,\mu,\delta)=(1,0,0.1)$, $\gamma T=0.025$, $c=0.2$, $R_0=1.38$.}
\label{p:fig:growth-phase}
\end{figure}

Note that the early lag is an eclipse phase that nobody put in. Within-host
models routinely insert one by hand, as a fitted delay distribution or as
explicit eclipse compartments \cite{ciupe2024incorporating,herz1996viral};
here it emerges from the intracellular dynamics with no extra parameter, its
scale set by $\mean{T_{\rm prod}}$. The caution that goes with the
observation is Herz et al.'s \cite{herz1996viral}: the shoulder in
post-treatment viral load is a compound of pharmacological delay,
intracellular delay, infected-cell decay and viral clearance, so the honest
claim is not that this lag \emph{is} the eclipse but that a lag of this
shape arises without being fitted.

\subsection{Identifiability: a negative result}
\label{p:sec:identifiability}

How much of the intracellular process can a population-level fit recover? At
any fixed growth rate $r$, the map
\begin{equation*}
(\lambda,\mu,\delta)\ \longmapsto\ \lrb{p_{\rm eff}(r),\,d_{\Icell,{\rm eff}}(r)}
\end{equation*}
sends three rates to two numbers, and is therefore not invertible
(\cref{p:fig:identifiability}). This is a structural identifiability
statement of the kind Miao et al.\ \cite{miao2011identifiability} treat for
viral dynamics generally, and it is negative for the same reason: the
observable is a projection, and the projection has a kernel.

\begin{figure}[tbp]
\centering
\includegraphics[width=0.94\textwidth]{N4b_2_identifiability_levels.pdf}
\caption{Population observables leave level sets rather than identifying the
intracellular rates. With $c=1$ and $\gamma T=1$, the interior flooding
point $(\lambda,\mu,\delta)=(1,0.5,1/3)$ has $R_0=2$ and $r_\star=0.294$.
(a) One population observable leaves a level curve: in the slice $\lambda=1$
the contours $R_0=2$ and $r=r_\star$ are distinct; alternatives A and C
share $R_0$ with the reference but not its growth rate. (b) With $\lambda$
free, a continuous family of triples shares both. Joint knowledge of
threshold and exponential growth therefore does not identify
$(\lambda,\mu,\delta)$.}
\label{p:fig:identifiability}
\end{figure}

It is worth seeing that this result is corroborated from the fitting side.
Mittler et al.\ \cite{mittler1998influence} compared a range of
intracellular delay distributions against HIV RNA data after protease
inhibitor initiation and found the resulting parameter estimates largely
insensitive to which distribution was assumed --- a finding influential
enough that the review of \cite{ciupe2024incorporating} cites it to justify
preferring the simplest eclipse compartment. Read through
\cref{p:thm:projection}, that insensitivity is not a rival claim but the
same one: on a single exponential phase a constant pair reproduces the
kernel exactly, so the shape of the kernel is invisible in exactly the
regime where those fits were performed. Consistently, when Kakizoe et al.\
\cite{kakizoe2015method} fitted in vitro SHIV data resolving the early
transient, an Erlang delay did improve the fit.

A third observable is therefore required; in order of experimental
accessibility:
\begin{enumerate}[leftmargin=1.9em]
\item the burst-size dispersion --- the geometric ratio $1/a$ fixes one
combination of the rates directly;
\item the fraction of cells that clear without bursting, which is $b$;
\item the shape --- not merely the mean --- of the release-time density
$\delta J$.
\end{enumerate}
Any one of these, together with $(p,d_{\Icell})$, over-determines
$(\lambda,\mu,\delta)$ and converts the fit into a consistency test of the
intracellular model. The limitation is not one of fitting procedure but of
the information content of standard BMVR itself, and the remedy is the one
\cite{ciupe2024incorporating} calls for in general terms: data at a finer
scale of resolution, chosen to match the model's structure.

\FloatBarrier

\section{Bursting versus budding: the flooding parameter~\texorpdfstring{$L$}{L}}
\label{p:sec:flooding}

Does bursting help a pathogen get established? The question belongs to the
stochastic early phase, where counts are small and extinction is a live
possibility \cite{pearson2011stochastic}, and it is the question the
stochastic intracellular model was built to answer: both versions of the
early infection are branching processes whose offspring distributions are in
closed form.

The mechanism at issue is worth naming before it is analysed. A bursting
cell delivers its whole output in one instant, so a founder that survives to
rupture presents the target population with many simultaneous infection
attempts rather than a trickle spread over its lifetime. Whether that
concentration helps or hurts is not obvious: many attempts at once buy
insurance against the loss of any one particle, but they are all forfeited
if the intracellular population dies out before the burst, which happens
with probability $b$. We call the concentration \emph{flooding}, and the
question is when it pays.

\subsection{The offspring of one infected cell}

In the branching approximation each released particle independently either
infects a target cell, at rate $\gamma T$, or is cleared, at rate $c$, and so
infects with probability
\begin{equation}
q=\frac{\gamma T}{\gamma T+c}.
\label{p:eq:q}
\end{equation}
The offspring of one infected cell is the burst size $\Kb$ thinned by $q$.
Summing \cref{p:eq:burstlaw}, the unconditional burst-size \pgf{} is
$G_{\Kb}(z)=b+(\delta/\lambda)z/(a-z)$ --- including the mass $b$ of cells
that never burst --- and thinning replaces $z$ by $y:=1-q+qz$:
\begin{equation}
G_{\rm off}(z)=b+\frac{\delta}{\lambda}\,\frac{y}{a-y},
\qquad y=1-q+qz.
\label{p:eq:Goff}
\end{equation}
One checks $G_{\rm off}(1)=b+(\delta/\lambda)/(a-1)=b+B=1$. The mean and
variance of the offspring number are
\begin{equation}
m=G_{\rm off}'(1)=q\,V_\infty,
\qquad
\mathrm{Var}_{\rm burst}
=\frac{2q^2V_\infty}{a-1}+qV_\infty-q^2V_\infty^2.
\label{p:eq:offmoments}
\end{equation}
When $\gamma T\ll c$, $q\approx\gamma T/c$ and $m\approx R_0$.

The comparator is the budding cell of \cref{p:sec:budding}, whose output is
geometric with mean $p/d_{\Icell}$. Match it to the bursting cell's total
output, $p/d_{\Icell}=V_\infty$. That matching is what makes the comparison
fair, and by \cref{p:thm:R0inv} it costs nothing: at equal total output the
two mechanisms have exactly the same $R_0$, so everything that follows
compares two pathogens at one threshold, and any difference between them is
a difference in how the output is \emph{scheduled} rather than in how much of
it there is. Thinning a geometric law leaves a geometric law, so the budding
offspring number is geometric with the same mean $m=qV_\infty$.

\subsection{Extinction, and the flooding advantage}

The extinction probability of a branching process is the smallest root of
$G_{\rm off}(z)=z$ in $[0,1]$. With $y=1-q+qz$ the fixed-point equation is
quadratic, $z=1$ is always a root, and the other root gives, for $m>1$,
\begin{equation}
z_{\rm ext}^{\rm burst}
=\frac{a(1-q+qb)-1+q}{q},
\qquad
z_{\rm ext}^{\rm bud}=\frac{1}{m}.
\label{p:eq:zext}
\end{equation}
Everything now turns on a single combination of the intracellular rates.

\begin{definition}[The flooding parameter]
\label{p:def:L}
\begin{equation}
L:=a(1-b)=a-\frac{\mu}{\lambda},
\qquad\text{so that}\qquad
V_\infty=\frac{L}{a-1}.
\label{p:eq:Ldef}
\end{equation}
\end{definition}

\begin{theorem}[The flooding advantage]
\label{p:thm:flood}
For $m>1$,
\begin{equation}
z_{\rm ext}^{\rm burst}-z_{\rm ext}^{\rm bud}
=\bigl(L-1\bigr)\Bigl(\frac{1}{m}-1\Bigr).
\label{p:eq:flood}
\end{equation}
Hence bursting strictly lowers the establishment extinction probability,
relative to matched-mean budding, if and only if $L>1$, that is
\begin{keyresult}
\begin{equation}
\delta(\lambda+\mu)>\lambda\mu
\qquad\Longleftrightarrow\qquad
\frac{1}{\delta}<\frac{1}{\lambda}+\frac{1}{\mu},
\label{p:eq:floodcrit}
\end{equation}
\end{keyresult}
with the convention that the condition holds for all $\delta>0$ when
$\mu=0$. When $L=1$ the two extinction probabilities are identical for every
$q$.
\end{theorem}

\begin{proof}
The identity \cref{p:eq:flood} is algebra: substitute \cref{p:eq:zext},
$m=qV_\infty=qL/(a-1)$ and simplify. Since $m>1$ makes the second factor
negative, the sign of the difference is the sign of $1-L$. For the
criterion, $L>1$ is $a(1-b)>1$, i.e.\ $a>1+\mu/\lambda$; inserting the
definition of $a$ from \cref{p:eq:roots} and squaring --- legitimate since
both sides are positive, or the right-hand side is not --- yields
$(\lambda+\mu+\delta)^2-4\lambda\mu>4\lambda^2(1+\mu/\lambda-\eta)^2$, which
reduces to $\delta(\lambda+\mu)>\lambda\mu$.
\end{proof}

The criterion \cref{p:eq:floodcrit} is the chapter's central biological
statement, and it reads cleanly as a comparison of timescales. The burst clock $1/\delta$ must
be shorter than the harmonic-type combination
$\bigl(1/\lambda+1/\mu\bigr)$ of the birth and death timescales. A cell that
takes too long to rupture is a cell whose internal population has time to
die out first, and the all-or-nothing hazard $b$ then outweighs the
insurance that a simultaneous batch buys. When $\mu=0$ that hazard vanishes
and bursting wins at every $\delta$. \Cref{p:fig:L-landscape} shows where
the boundary lies in the rate plane; \cref{p:tab:flooding} confirms the
identity in each regime, and \cref{p:fig:flooding-regimes} traces both
extinction probabilities across the whole range of $q$.

\begin{figure}[tbp]
\centering
\includegraphics[width=0.64\textwidth]{N4b_4_L_landscape.pdf}
\caption{Where a pathogen sits relative to the flooding boundary. The
flooding parameter $L(\mu,\delta)$ at $\lambda=1$; the exact curve $L=1$,
equivalently $\delta_*=\lambda\mu/(\lambda+\mu)$, partitions the rate plane
into a region where bursting lowers extinction risk and one where it raises
it. The three marked points are the regimes of \cref{p:tab:trio},
$(\mu,\delta)=(0,0.1)$, $(0.5,1/3)$ and $(0.9,0.1)$, with $L=1.10$, $1.00$
and $0.42$.}
\label{p:fig:L-landscape}
\end{figure}

\begin{table}[htbp]
\centering
\renewcommand{\arraystretch}{1.1}
\begin{tabular}{@{}l S[table-format=1.1]
    S[table-format=1.3]S[table-format=1.3] l@{}}
\toprule
{$(\lambda,\mu,\delta)$} & {$q$} & {$z_{\rm ext}^{\rm burst}$}
 & {$z_{\rm ext}^{\rm bud}$} & {verdict} \\
\midrule
$(1,0,0.1)$   & 0.2 & 0.400 & 0.455 & bursting wins ($L=1.10$) \\
$(1,0.9,0.1)$ & 0.9 & 0.935 & 0.844 & budding wins ($L=0.42$) \\
$(1,0.5,1/3)$ & 0.6 & 0.833 & 0.833 & tie ($L=1$) \\
$(1,0.5,1/3)$ & 0.9 & 0.556 & 0.556 & tie ($L=1$) \\
\bottomrule
\end{tabular}
\caption{Extinction probabilities from \cref{p:eq:zext} and the sign of
their difference. The gap is $(L-1)(1/m-1)$ in every row.}
\label{p:tab:flooding}
\end{table}

\begin{figure}[tbp]
\centering
\includegraphics[width=0.94\textwidth]{F4b_2_flooding_regimes.pdf}
\caption{Extinction probabilities for bursting and matched-mean budding
across the three regimes of \cref{p:tab:trio}. Each pair of curves begins at
$m=1$; the signed gap follows \cref{p:eq:flood}. The $q$-windows differ
because $m=1$ sits at different $q$ in each regime.}
\label{p:fig:flooding-regimes}
\end{figure}

For $m\le1$ both processes die out almost surely and the comparison is
empty; \cref{p:eq:flood} changes sign there, an artefact of both formulae
running outside their domain $[0,1)$, not a reversal of the biology.

\subsection{The boundary \texorpdfstring{$L=1$}{L=1}: a coincidence of laws}

The boundary case is more than a tie of extinction probabilities.

\begin{proposition}
\label{p:prop:L1}
The following are equivalent:
\begin{equation}
L=1
\;\Longleftrightarrow\;
a=1+\frac{\mu}{\lambda}
\;\Longleftrightarrow\;
\delta(\lambda+\mu)=\lambda\mu
\;\Longleftrightarrow\;
\frac{1}{\delta}=\frac{1}{\lambda}+\frac{1}{\mu}.
\label{p:eq:L1equiv}
\end{equation}
When they hold, the bursting offspring law \cref{p:eq:Goff} is
\emph{exactly geometric} with mean $m$ --- not merely equal in extinction
probability to the budding comparator. In particular the offspring variances
coincide, and in general
\begin{equation}
\mathrm{Var}_{\rm bud}-\mathrm{Var}_{\rm burst}
=\frac{2q^2V_\infty(L-1)}{a-1},
\label{p:eq:varorder}
\end{equation}
so that at matched mean the bursting offspring variance is smaller than the
budding one when $L>1$, equal when $L=1$, and larger when $L<1$.
\end{proposition}

\begin{proof}
$L=1$ is $a(1-b)=1$, i.e.\ $a-1=ab=\mu/\lambda$, giving the chain of
equivalences. Under it $b=\delta/\lambda$, and \cref{p:eq:Goff} collapses:
\[
G_{\rm off}(z)=b\Bigl(1+\frac{y}{a-y}\Bigr)=\frac{ab}{a-y}
=\frac{a-1}{a-y}
=\frac{(a-1)/a}{1-y/a},
\]
using $ab=a-1$. With $y=1-q+qz$,
$1-y/a=\frac{a-1+q}{a}\bigl(1-\frac{q}{a-1+q}z\bigr)$, so
\[
G_{\rm off}(z)=\frac{1-\varrho}{1-\varrho z},
\qquad
\varrho=\frac{q}{a-1+q}=\frac{m}{1+m},
\qquad m=\frac{q}{a-1}=qV_\infty,
\]
the \pgf{} of a geometric law on $\{0,1,2,\ldots\}$ with mean $m$ --- the
same law as the matched budding comparator. The variance identity
\cref{p:eq:varorder} is the difference between the geometric variance
$m(1+m)$ and \cref{p:eq:offmoments}, with $V_\infty=L/(a-1)$.
\end{proof}

\noindent
So $L$ does more than fix a sign. At the boundary, bursting and budding are
the same branching process in different clothes; away from it,
\cref{p:eq:varorder} says precisely how the two offspring laws differ, and
the extinction ordering of \cref{p:thm:flood} is the ordering of their
variances at fixed mean. That is the mechanism behind the criterion:
concentrating output into one batch reduces offspring variance --- which
helps at establishment --- but only if the batch is delivered often enough
to offset the mass $b$ at zero.

\subsection{The price: invasion speed}
\label{p:sec:tradeoff}

Establishment is half the story; the other half is how fast an established
infection grows. Match the models a second time, at $R_0$ and at mean
productive lifetime: take the budding cell with
$d_{\Icell}=1/\mean{T_{\rm prod}}$ and $p=V_\infty d_{\Icell}$ --- exactly
$p_{\rm eff}(0)$, the fairest constant match available --- and choose
$\gamma T$, or $c$, so that both share $R_0=2$. Solving for the growth rates,
$r_{\rm bud}$ from $(r+d_{\Icell})(r+c)=\gamma Tp$ and $r_{\rm burst}$ from
\cref{p:eq:char}, gives the same ordering in all three regimes:
\begin{equation}
r_{\rm bud}>r_{\rm burst}.
\label{p:eq:rorder}
\end{equation}

\begin{table}[htbp]
\centering
\renewcommand{\arraystretch}{1.1}
\begin{tabular}{@{}l S[table-format=1.2]
    S[table-format=1.3]S[table-format=1.3]S[table-format=1.3]@{}}
\toprule
{$(\lambda,\mu,\delta)$} & {$L$} & {$d_{\Icell}$}
 & {$r_{\rm bud}$} & {$r_{\rm burst}$} \\
\midrule
$(1,0,0.1)$   & 1.10 & 0.417 & 0.250 & 0.180 \\
$(1,0.5,1/3)$ & 1.00 & 0.910 & 0.395 & 0.294 \\
$(1,0.9,0.1)$ & 0.42 & 0.701 & 0.343 & 0.198 \\
\bottomrule
\end{tabular}
\caption{Deterministic growth rates at matched $R_0=2$, matched total yield
and matched mean productive lifetime, with $c=1$. Budding invades faster on
both sides of the flooding boundary and on it, so the ordering is not the
same effect as \cref{p:thm:flood}.}
\label{p:tab:tradeoff}
\end{table}

Whether \cref{p:eq:rorder} holds in general is not settled here, but it
reduces to a single inequality in which the coupling constants do not appear.

\begin{proposition}[The ordering is a Laplace-transform comparison]
\label{p:prop:rorder}
Match the comparator as above, so that
$\int_0^\infty g_{\rm bud}=\int_0^\infty g_{\rm burst}=V_\infty$. Then
$r_{\rm bud}>r_{\rm burst}$ for \emph{every} pair $(\gamma T,c)$ with
$R_0>1$ if and only if
\begin{equation}
\Lap{g}_{\rm burst}(r)\ \le\ \Lap{g}_{\rm bud}(r)=\frac{p}{r+d_{\Icell}}
\qquad\text{for all }r>0,
\label{p:eq:laplace-order}
\end{equation}
that is, if and only if the release age under bursting dominates the matched
exponential release age in the Laplace-transform order.
\end{proposition}

\begin{proof}
Each model's growth rate is the unique root of
$\phi(r):=\gamma T\Lap g(r)-(r+c)$; $\phi$ is strictly decreasing, and
$\phi(0)=\gamma TV_\infty-c=c(R_0-1)>0$, so the root is positive and unique.
If \cref{p:eq:laplace-order} holds then $\phi_{\rm burst}\le\phi_{\rm bud}$
pointwise, so the root of the former is no larger. Conversely, suppose
$\Lap g_{\rm burst}(r_0)>\Lap g_{\rm bud}(r_0)$ for some $r_0>0$. Choose any
$c>0$ and set $\gamma T=(r_0+c)/\Lap g_{\rm burst}(r_0)$, so that
$r_{\rm burst}=r_0$; then $\phi_{\rm bud}(r_0)<0$ and $r_{\rm bud}<r_0$.
Taking $c$ small makes $R_0=\gamma TV_\infty/c$ exceed one, so the ordering
fails for an admissible coupling.
\end{proof}

\noindent
Both sides of \cref{p:eq:laplace-order} agree at $r=0$, by the matching, so
its content is entirely in how the two release schedules are spread over
age. The inequality was checked at $1500$ triples over
$(\mu,\delta)\in[10^{-3},10]\times[10^{-3},16]$ at $\lambda=1$, at $60$
values of $r$ across four decades, without a violation
(test~N of the verification suite). The \emph{usual} stochastic order does
not hold --- the burst release-age survival function dips below the matched
exponential --- so \cref{p:eq:laplace-order} is not a weaker way of saying
something stronger. It is the order the problem actually has.

The mechanism is generation time in the exact sense of
\cref{p:sec:generation-kernel}. The budding cell starts releasing at once
while the bursting cell must first accumulate load, so budding's
generation-time density $\mathcal A/R_0$ sits earlier, and by
$\Lap{\mathcal A}(r)=1$ an earlier density gives a larger root. At
$(\lambda,\mu,\delta)=(1,0,0.1)$ with $c=1$ the mean generation time is
$2.398+1=3.398$ for budding against $3.089+1=4.089$ for bursting, which is
what puts $r_{\rm bud}=0.250$ above $r_{\rm burst}=0.180$
(\cref{p:fig:generation-times}). This is the same delay that appeared as the
eclipse-like lag of \cref{p:sec:overlays}, and it is the phage lysis-timing
problem in another notation: holding on longer makes more particles per cell
but takes longer to make them \cite{wang2006lysis,wallinga2007generation}.
That an infectious-period distribution shifts the growth rate at fixed
threshold is a general fact about renewal epidemics rather than anything
peculiar to bursting \cite{lloyd2001realistic,champredon2015intrinsic}; what
is particular here is that the distribution is derived rather than assumed.

\begin{figure}[tbp]
\centering
\includegraphics[width=0.94\textwidth]{N4b_6_generation_times.pdf}
\caption{The generation-time structure behind the first row of
\cref{p:tab:tradeoff}. At $(\lambda,\mu,\delta)=(1,0,0.1)$, budding and
bursting are matched at $R_0=2$, $c=1$, total yield and mean productive
lifetime. The main curves are the normalised infection-generation laws
$\mathcal A(\alpha)/R_0$, including the common extracellular waiting-time
component; panel (b) isolates release age and shows why the budding law
leads.}
\label{p:fig:generation-times}
\end{figure}

\subsection{The trade-off}

Put together, the two comparisons say what bursting buys and what it costs.
At fixed total output per cell, bursting leaves $R_0$ untouched, lowers the
probability of early extinction when $L>1$, and slows deterministic
invasion. A pathogen tuning $\delta$ therefore faces a genuine trade-off
(\cref{p:fig:pareto}): $z_{\rm ext}$ and $r$ move in opposite directions,
and no setting of $\delta$ improves both. For $\mu>0$, budding beats
bursting on extinction until $\delta$ crosses
$\delta_*=\lambda\mu/(\lambda+\mu)$, beyond which bursting takes over; for
$\mu=0$ bursting wins at every $\delta$, at the continuing cost of a slower
$r$.

\begin{figure}[tbp]
\centering
\includegraphics[width=0.62\textwidth]{N4b_5_pareto_extinction_growth.pdf}
\caption{Tuning $\delta$ traces a strict extinction--growth trade-off, drawn
parametrically as $\delta$ varies from $0.03$ to $0.8$ with $\lambda=1$,
$\mu=0$, $c=1$ and $R_0=2$. The coupling is reset along the curve to hold
$R_0=2$, and the budding comparator is matched in lifetime yield $V_\infty$
and in mean productive lifetime. There is no setting of $\delta$ that
improves both coordinates: faster growth is bought with a higher extinction
probability.}
\label{p:fig:pareto}
\end{figure}

Two remarks place the result. First, its scope: both comparisons hold the
offspring \emph{mean} fixed and are branching-process statements, so they
say nothing about nonlinear effects, where other advantages of flooding may
live (\cref{p:sec:limitations}). Second, its company. Lord and Bonsall
\cite{lord2021evolutionary} reach a qualitatively parallel conclusion for
within-host virus evolution by a different route, using fixed-delay models
to compare virion release strategies and finding that lysis-like strategies
are favoured when the terminal event comes quickly and budding when it is
slow. The criterion \cref{p:eq:floodcrit} is a sharp version of
``quickly'', written in intracellular rates, and it comes with the exact size of the advantage from
\cref{p:eq:flood}. The evolutionary framing belongs to \ChPath{}, which
inherits from this comparison the awkward fact that one mechanism can be
favoured on establishment and disfavoured on invasion speed at once --- the
nested-model problem in its usual form \cite{gilchrist2006evolution}.

\FloatBarrier

\section{The reach of the construction}
\label{p:sec:reach}

\Cref{p:sec:generality} showed that any pair of kernels induces the same
population equations. Bursting and constant budding are the ends of a
spectrum; between them lie mechanisms that release many times per cell or
release part of the load, and beyond it lie models that are not
birth--death--catastrophe processes at all. All supply a pair $(S,g)$, and
none changes \cref{p:eq:skeleton}. This section works through the two
intermediates that matter and then says where the generality stops.

\subsection{Recurrent release: the reset process}
\label{p:sec:reset}

The absorbing process releases at most once per cell. The intermediate worth
isolating lets a cell release \emph{many} times: while the cell is alive the
load undergoes birth $\lambda n$, death $\mu n$, optional immigration $\nu$,
and a \emph{dump that resets} --- at rate $\delta n$ the whole pool is
released and the load returns to zero, after which it may re-accumulate ---
with the cell dying at a separate, load-independent rate $d_{\Icell}$.
Reading the load as a local assembled pool rather than the whole-cell
contents makes this a tractable multi-pulse caricature between full dumping
and continuous export.

Because dumps do not absorb, the kernels change but the skeleton does not:
\begin{equation}
S(\alpha)=\ee^{-d_{\Icell}\alpha},
\qquad
g(\alpha)=\delta\,\ee^{-d_{\Icell}\alpha}\,\mean{X_\alpha^2},
\label{p:eq:reset-kernels}
\end{equation}
with $\mean{X_\alpha^2}$ computed for the non-absorbing reset process. Linear
rates keep its moments tractable, so $g$ is explicit once they are solved,
and structurally $g\propto\mean{X^2}$ exactly as for the absorbing process.
The effective parameters become
\begin{equation}
p_{\rm eff}(r)=\lrb{r+d_{\Icell}}\Lap{g}(r),
\qquad
d_{\Icell,{\rm eff}}\equiv d_{\Icell},
\label{p:eq:reset-peff}
\end{equation}
and $R_0=(\gamma T/c)\int_0^\infty g$, finite only for $d_{\Icell}>0$ since
otherwise recurrent dumps could continue forever. If the reset process
admits a stationary law $\pi$ and $d_{\Icell}$ is small, the mature-cell
limit is
\begin{equation}
p_\infty=\lim_{\alpha\to\infty}\frac{g(\alpha)}{S(\alpha)}
=\delta\,\mathbb{E}_\pi[X^2],
\label{p:eq:reset-mature}
\end{equation}
the same form as the old-cell limit \cref{p:eq:oldcell} of the absorbing
process. A long-lived infected cell therefore behaves as constant-$p$ BMVR
with that $p$, while transients --- first dumps, young cells --- still need
the full kernel. This is what makes the spectrum claim concrete rather than
gestural: an intermediate with kernels, an effective release rate, a mature
limit, and a definite answer as to which standard model it converges on.

\subsection{Partial release, and antiviral mode of action}
\label{p:sec:partial}

The remaining axis is \emph{how much} leaves the cell when a release event
fires. The natural intermediate removes a random part of the load; the
default is boolean thinning, each unit released independently with
probability $\varphi$. At the mean-field level with event clock $\delta X$,
writing $Q$ for the total intracellular stock across productive cells, the
population equations are
\begin{equation}
\ddt{\Icell}=\gamma T\,\Vfree-d_{\Icell}\Icell,
\qquad
\ddt{Q}=\nu\Icell-\delta\varphi\,\frac{Q^2}{\Icell}-d_{\Icell}Q,
\qquad
\ddt{\Vfree}=\delta\varphi\,\frac{Q^2}{\Icell}-c\,\Vfree.
\label{p:eq:partial}
\end{equation}
With a load-independent clock $\lambda(X)=\delta$ instead, the mean export is
$\delta\varphi Q$ and the equations collapse to stock-and-export dynamics
with $\varepsilon_{\rm eff}=\delta\varphi$. So partial release is invisible
in the mean-field equations whenever the event clock is load-independent,
and visible when it is load-dependent; in single-cell release-count
distributions it is always visible.

This axis is where the construction meets the second of the two payoffs that
\cite{ciupe2024incorporating} identifies for intracellular modelling: the
ability to locate a drug's mode of action at a named step of the viral life
cycle. In the multiscale hepatitis~C models
\cite{guedj2013modeling,rong2013mathematical} a direct-acting antiviral is
represented by three separate factors --- blocking intracellular RNA
production by $(1-\epsilon_a)$, blocking packaging and secretion by
$(1-\epsilon_s)$, and enhancing intracellular degradation by a factor at
least one. Each has a counterpart here. Blocking production lowers
$\lambda$; enhancing degradation raises $\mu$; blocking entry lowers
$\gamma T$; and blocking secretion is exactly a reduction in $\varphi$,
since $\varphi=1-\epsilon_s$. Every one of these is a transformation of
$(S,g)$, so a drug in this framework does not add parameters to the
population layer but moves the kernels, and its population signature is
whatever \cref{p:eq:peff} makes of the new pair.

Two consequences follow, both in principle testable. The first is that a
drug moves the pathogen relative to the flooding boundary, and the direction
is read off \cref{p:eq:floodcrit} without further calculation. Enhancing
intracellular degradation raises $\mu$ and so lowers
$1/\lambda+1/\mu$, moving the pathogen toward the $L<1$ side on which its
burst schedule is a liability rather than an asset at establishment;
blocking intracellular production raises $1/\lambda$ and moves it the other
way. Neither move need cross the boundary --- \cref{p:eq:floodcrit} is an
inequality, not a monotone score --- but the sign of the shift is fixed, and
\cref{p:eq:flood} gives the size of what is gained or lost at each end. The
second is that because $p_{\rm eff}$ depends on $r$, an efficacy inferred
from a decay slope inherits the same phase dependence that
\cref{p:sec:span} finds in $p$ itself, which is one reading of why the
hepatitis~C clearance estimates of \cref{p:sec:fitting-consequences}
disagreed across therapies. Whether the flooding advantage survives at
intermediate $\varphi$, and whether the closed forms persist there at all,
is open.

\subsection{Beyond birth--death--catastrophe: HIV}
\label{p:sec:hiv}

The results above are written for pathogens whose intracellular life ends in
load-proportional catastrophe. HIV is not one: it buds from the plasma
membrane, and an infected cell need not rupture in a single event. After
integration the template count is one or a few proviruses and structural
products do not self-replicate, so the mean load follows an
immigration-like $\bar x'\approx\nu-\mu\bar x$ rather than the multiplicative
growth that produces $\delta K$. Yet limiting-dilution data show something
surge-like --- long quiet periods, then a concentrated bout of release ---
which is the phenomenon the catastrophe process exists to explain. Hataye
et al.\ \cite{hataye2019principles} report first detection around day~3 with
bulk production over days~3--5, an Erlang eclipse with $n\approx5$ stages,
many silent lineages, and a sigmoid rather than linear dependence of
establishment on the number of virus-releasing cells, with a threshold near
$2.3$ cells.

The position taken here is neither to force HIV into the catastrophe process
nor to retreat to constant-$p$ BMVR. The skeleton \cref{p:eq:skeleton}
carries HIV on a different pair of kernels: $S_{\rm HIV}$ the probability of
still being in a productive lineage and $g_{\rm HIV}$ the density of
productive sojourn times, with a multi-stage eclipse supplying the delay ---
near zero for small $\alpha$, then rising and falling. The effective
parameters carry over verbatim under
$(\Ifix,g)\mapsto(S_{\rm HIV},g_{\rm HIV})$. What does not carry over is
independent branching near the observed threshold: the data favour a sigmoid
incidence such as
\begin{equation}
f(\Vfree,T)=\gamma T\,\Vfree\cdot\frac{\Vfree}{K_A+\Vfree},
\label{p:eq:allee}
\end{equation}
inefficient at very low $\Vfree$ and saturating to mass action for
$\Vfree\gg K_A$. It is worth noting that this is the opposite curvature at
low $\Vfree$ from the saturating alternatives catalogued in
\cite{ciupe2024incorporating} --- Holling type~2, Beddington--DeAngelis,
Crowley--Martin --- which depress infection at \emph{high} viral load; the
departure from mass action wanted here is an Allee effect at establishment,
not saturation at saturation. Under \cref{p:eq:allee}, $R_0$ remains the
large-population threshold while early establishment is governed in addition
by the scale $K_A$, and \cref{p:prop:first-moment} no longer holds exactly,
since a nonlinear incidence breaks the moment closure.

\subsection{Where the reach stops}

What carries over to any single-cell mechanism is \cref{p:eq:skeleton}, the
map $p_{\rm eff}=\Lap g/\Lap S$, the reading of constant $p$ as an
exponential-phase projection, and the generation-kernel formulation of $R_0$
and $r$. What does not is everything resting on the closed forms of the
birth--death--catastrophe process: the flooding criterion
$1/\delta<1/\lambda+1/\mu$, the geometric burst law and its identification
with the quasi-stationary distribution, and one-shot cumulative release as
the only single-cell trajectory. For HIV in particular, independent
branching is not the full establishment story near the observed threshold;
the flooding comparison survives as a \emph{linear} baseline, at matched
mean yield with offspring variances compared, but synergy can dominate
variance effects there and the two mechanisms must be separated before
either is invoked.

\section{Discussion}
\label{p:sec:discussion}

This chapter took a completed stochastic single-cell theory and embedded it
in a population. What came out is a renewal model of within-host infection
forced by the intracellular process and exact in its kernels; the
demonstration that the standard model is a projection of it, valid one
exponential phase at a time; an effective-parameter map whose two ends
differ by the factor $\mean{X^2}_{\rm QS}$; and a closed-form answer to
whether bursting pays, in which the flooding parameter $L=a(1-b)$ decides
the sign and $(L-1)(1/m-1)$ gives the size.

\subsection{Consequences for fitting practice}
\label{p:sec:fitting-consequences}

Three results bear on how within-host models are fitted, and each says what
to do rather than only what is true.

\paragraph{A fitted $p$ is a property of the phase.} Acute-phase and
set-point estimates should therefore differ systematically, in a direction
and by an amount \cref{p:eq:peff} specifies. Two published estimates of $p$
for the same organism at different phases are not in conflict, and the ratio
between them is itself an observable: by \cref{p:eq:peff-span} it cannot
exceed $\mean{X^2}_{\rm QS}$, so a reported spread larger than that is
evidence against the intracellular model rather than against either fit.
There is precedent for the general shape of this argument. Hepatitis~C
clearance was estimated at $c=6.2$/day from interferon-$\alpha$ trials
\cite{neumann1998hepatitis} and at $c=23.3$/day from direct-acting antiviral
data \cite{guedj2013modeling} --- close to a fourfold spread in a rate that
cannot depend on the therapy used to measure it. The resolution was not to
adjudicate between the fits but to add intracellular structure, after which
the two biphasic declines turn out to be reporting different pairs of
underlying rates: under interferon-$\alpha$ the slopes are the intracellular
loss rate and the infected-cell death rate, and neither is $c$, whereas
under direct-acting antivirals the fast slope is $c$ after all
\cite{rong2013mathematical}. The earlier figure was not a bad measurement of
clearance; it was a good measurement of something else. That is the same
error \cref{p:sec:projection} identifies for $p$, and it is the general
lesson: a slope read off a transient is not a rate of the organism until a
model says which rate it is.

\paragraph{The eclipse may be being fitted twice.} The lag that within-host
models insert by hand emerges here from the intracellular dynamics with no
extra parameter, its scale set by $\mean{T_{\rm prod}}$. A model already
carrying a fitted eclipse and a bursting pathogen is therefore fitting the
same delay twice, and the two estimates should be checked against each other
before either is trusted.

\paragraph{A population fit needs a single-cell companion.} No fit of
$(p,d_{\Icell})$ recovers the three intracellular rates
(\cref{p:sec:identifiability}), so a population fit is worth doing only
alongside one of the three single-cell observables named there. With any one
of them the system is over-determined and the fit becomes a test of the
intracellular model rather than a description of it. This is a concrete
instance of the general call in \cite{ciupe2024incorporating} for data
collected at a finer scale of resolution to validate complex models: the
chapter names which three measurements would do it, and in what order of
experimental accessibility.

\subsection{Limitations}
\label{p:sec:limitations}

The assumptions of \cref{p:sec:assumptions} are load-bearing. Holding target
cells constant (A1) makes the model linear, so supercritical solutions grow
without bound and peak times exist only in subcritical cases; depletion
would cap them, but the linear analysis remains the right object for
establishment and early growth, where the chapter's claims live. The
continuum convention (A2) is why $R_0=\gamma TV_\infty/c$ sits alongside the
branching mean $m=qV_\infty$: they agree when $\gamma T\ll c$ and are
otherwise different quantities, and confusing them would misread
\cref{p:thm:flood}. The synchronous initial condition (A3) is not innocuous,
since a non-trivial initial age distribution adds a term to
\cref{p:eq:Icell}.

Single infection (A4) is the most substantive restriction. Multiply infected
cells release superlinearly more --- \cref{p:eq:Vk} gives $V_\infty^{(k)}$
--- and multiplicity is not rare: as many as a tenth of HIV-infected cells
are thought to be multiply infected, and coinfection shapes estimates of
infected-cell death and production, drug-resistance evolution and immune
escape \cite{ciupe2024incorporating}. The standard remedy adds compartments
for doubly and triply infected cells, which the same review notes is
difficult to validate against plasma titres because of the parameter count.
The MOI-$k$ row of \cref{p:tab:sg} suggests a cheaper route: multiplicity
enters through the kernels
$\bigl(I^k-D^k,\ \delta K_k\bigr)$ while the skeleton stays
one-dimensional, adding no compartments and no fitted parameters. Carrying
those kernels through the population layer has not been done here, and it is
the most obvious extension. It would also connect to the alternative
low-dimensional treatment of coinfection due to Koelle et al.\
\cite{koelle2019within}, whose macroparasite-style bookkeeping tracks the
total internalised virus across all cells --- the quantity written $Q$ in
\cref{p:eq:partial} --- and which likewise abandons the standard assumption
that burst size is independent of the extent of coinfection. Absence of
spatial structure (A5) is the most consequential omission for a pathogen
that spreads cell to cell.

Finally, the chapter is written in dimensionless rates with $\lambda=1$ and
does not say where a real pathogen sits on \cref{p:fig:L-landscape}. For
\yp{} in the macrophage the intracellular interval is short and the biology
resists a clean measurement: the bacterium enters a permissive vacuole and
replicates there while the temperature shift from the flea to the host
induces the capsule, the pH~6 antigen and the type~III secretion system, and
then leaves by a route the primary literature still treats as unresolved,
with apoptosis, pyroptosis and necroptosis all reported. Placing that
organism relative to $L=1$ needs $\lambda$, $\mu$ and $\delta$ for the
intracellular phase --- and the exit mechanism is exactly what fixes
$\delta$.

\subsection{Open problems}
\label{p:sec:open-problems}

\begin{enumerate}[leftmargin=1.9em]
\item \textbf{Two-type killed second moments.} The release kernel of a
two-type renewal model needs
\begin{equation*}
\mean{X_\alpha^2\ind\{\tau_c>\alpha\}},\quad
\mean{X_\alpha Y_\alpha\ind\{\tau_c>\alpha\}},\quad
\mean{Y_\alpha^2\ind\{\tau_c>\alpha\}},
\end{equation*}
together with the two-type analogue of $D$. The moment hierarchy does not close upward, exactly as in
one type; the one-type escape was to differentiate the Riccati equation, and
the two-type route should be to differentiate the backward system of
\ChTwoType{}. This is the main technical obstacle between this chapter and
an eclipse-aware renewal model, and it should not be assumed easy.
\item \textbf{The flooding boundary against real parameters.} Locate \yp{}
in the macrophage relative to $\delta(\lambda+\mu)>\lambda\mu$ and compute
the size of the advantage from \cref{p:eq:flood}. The estimates discussed in
\ChDist{} --- the near-certain bursting of \ft{}
\cite{carruthers2020stochastic}, the clearance-dominated anthrax case ---
suggest real pathogens sit on both sides of the boundary, which is what
makes the question worth asking.
\item \textbf{Population-level variance.} Propagate $\mathrm{Var}(W)$
through the branching structure to a whole-host prediction. For single-cell
assays the first moments suffice; for whole-host data the variance is the
main stochastic signal that standard BMVR cannot reproduce, and the
integer-valued nature of the offspring count enters through the same
quantity \cite{williams2024reproduction}.
\item \textbf{Multiplicity through the population layer}, and
\textbf{partial release}: where in $\varphi$ the flooding advantage
survives.
\end{enumerate}

Two smaller questions are single definite statements rather than open
directions. The hypothesis of \cref{p:prop:peff-monotone} --- that
$\rho=g/S$ increases with cell age --- is a claim about the one closed-form
function \cref{p:eq:rho-closed} of one variable, so proving it is finite
algebra; and the ordering $r_{\rm bud}>r_{\rm burst}$ is by
\cref{p:prop:rorder} equivalent to the single Laplace inequality
\cref{p:eq:laplace-order}, with the coupling constants eliminated. Both are
supported by wide numerical sweeps and neither is proved.

\subsection{Connections forward}

The two-type process of \ChTwoType{} supplies the most immediate
continuation, and it is where the principal technical obstacle lives. In the
regime $\delta_1=0<\delta_2$ --- no burst before conversion --- the
no-catastrophe probability has zero initial hazard and then rises: precisely
the eclipse-phase signature that within-host models insert by hand, emerging
mechanistically, with the eclipse duration set by the conversion rate.
Nothing in \cref{p:eq:skeleton} need change to accept those kernels when
they exist. The bursting--budding comparison finds its evolutionary home in
\ChPath{}, where the intermediate mechanisms of \cref{p:sec:reach} supply
the phenotypes a selection argument needs. The nonlinear effects that
\fwd{nonlinear}{the later modelling chapters} treat --- immune saturation,
local depletion, the all-run-out-at-once intuition --- are exactly what the
branching comparison deliberately excludes, and should be named wherever the
flooding advantage is invoked.

\FloatBarrier

% ===========================================================================
%  Appendices
% ===========================================================================
\begingroup
\renewcommand{\theHsection}{\thechapter.appendix.\Alph{section}}
\renewcommand{\theHsubsection}{\theHsection.\arabic{subsection}}
\renewcommand{\thesection}{\thechapter.\Alph{section}}
\setcounter{section}{0}

\section{Single-cell results quoted from earlier chapters}
\label{p:app:quoted}

Write $w=\ee^{\theta t}$. The roots of $\lambda f^2-(\lambda+\mu+\delta)f+\mu$
are $a,b$ of \cref{p:eq:roots}, with $b<1<a$, and the derived constants are
\begin{equation}
A=a-1,\qquad B=1-b,\qquad \theta=\lambda(a-b),\qquad
\kappa=1+\frac{\delta}{2\lambda}.
\label{p:eq:derived}
\end{equation}
The identities that simplify everything downstream are
\begin{equation}
ab=\frac{\mu}{\lambda},\quad a+b=\frac{\lambda+\mu+\delta}{\lambda},\quad
AB=\frac{\delta}{\lambda},\quad A+B=a-b.
\label{p:eq:AB}
\end{equation}
The no-burst probability, the internal-extinction probability and the
non-fixation probability are
\begin{equation}
\begin{gathered}
I(t)=\frac{aB+bAw}{B+Aw},\qquad
D(t)=\frac{ab\,(w-1)}{aw-b},\\[0.35em]
\Ifix(t)=I(t)-D(t)=\frac{(a-b)^2\,w}{(B+Aw)(aw-b)}.
\end{gathered}
\label{p:eq:IDIhat}
\end{equation}
Here $I(\infty)=D(\infty)=b$, the probability that the cell clears
internally without ever bursting, while $\Ifix(0)=1$,
$\Ifix'(0)=-(\mu+\delta)$ and $\Ifix(\infty)=0$. With $I'=-\delta J$ and
$V'=\delta K$,
\begin{equation}
J(t)=\frac{(a-b)^2\,w}{(B+Aw)^2}=-\frac{\lambda}{\delta}(I-a)(I-b),
\label{p:eq:J}
\end{equation}
\begin{equation}
K(t)=\Bigl[1+\frac{2\lambda}{\delta}(1-I)\Bigr]J
=\frac{2\lambda}{\delta}(\kappa-I)J,
\label{p:eq:K}
\end{equation}
and the cumulative mean release is a quadratic in $1-I$ alone,
\begin{equation}
\begin{gathered}
V(t)=(1-I)\Bigl[1+\frac{\lambda}{\delta}(1-I)\Bigr],\\[0.35em]
V_\infty=\frac{a(1-b)}{a-1}=\frac{aB}{A},
\qquad
\mean{\Kb\mid\text{burst}}=\frac{a}{a-1},
\end{gathered}
\label{p:eq:V}
\end{equation}
in which $V_\infty$ counts the internally cleared cells with their release
of zero and the conditional mean removes them. The integral of the survival
kernel, which anchors \cref{p:sec:projection}, is
\begin{equation}
\mean{T_{\rm prod}}=\int_0^\infty\Ifix(\alpha)\,\dd\alpha
=\frac{1}{\lambda}\log\frac{a}{a-1}.
\label{p:eq:lifetime}
\end{equation}

\paragraph{Distribution-level facts.} The burst-time density is
$-I'(t)=\delta J(t)$, defective with total mass $1-b$. The burst size,
unconditional, obeys
\begin{equation}
\pr{\Kb=k}=\frac{\delta}{\lambda}\,a^{-k},\qquad k\ge1,
\label{p:eq:burstlaw}
\end{equation}
with the mass $b$ at ``no burst''; conditioned on bursting this becomes
$(a-1)a^{-k}$, exactly the quasi-stationary distribution of the load
\cite{van2011quasi} --- both geometric on $\{1,2,\ldots\}$ with ratio $1/a$.
Given that the burst occurs at time $t$, its size is a size-biased draw from
the state probabilities, $\mean{\Kb\mid\tau=t}=K(t)/J(t)$. The second moment
of the release, used for the offspring variance \cref{p:eq:offmoments}, is
\begin{equation}
\mean{W_t^2}
=\frac{2(\lambda-\mu)}{\delta}V-\frac{\lambda+\mu}{\delta}I-K
+\frac{\lambda+\mu}{\delta}+1.
\label{p:eq:EW2}
\end{equation}

\paragraph{Multiplicity of infection.} A cell founded by $k$ particles
releases superlinearly more:
\begin{equation}
g_k(t)=\delta K_k(t)
=\delta\lrb{kK I^{k-1}+k(k-1)J^2I^{k-2}},
\label{p:eq:gk}
\end{equation}
with lifetime yield, for $k\ge2$,
\begin{equation}
V_\infty^{(k)}
=(1-b^k)
+\frac{2\lambda}{\delta}\Biggl[\frac{1}{k+1}-b^k+\frac{k\,b^{k+1}}{k+1}\Biggr]
+k(k-1)\frac{\lambda}{\delta}\int_1^b (x-a)(x-b)\,x^{k-2}\,\dd x,
\label{p:eq:Vk}
\end{equation}
simplifying to $V_\infty^{(k)}=k+\lambda/\delta$ when $\mu=0$. The
conditional mean burst $V_\infty^{(k)}/(1-b^k)$ increases with $k$: multiply
infected cells burst more often, and bigger. The pair
$\bigl(I(\alpha)^k-D(\alpha)^k,\ \delta K_k(\alpha)\bigr)$ is the MOI-$k$
row of \cref{p:tab:sg}; assumption~(A4) sets it aside for the population
layer.

\section{Laplace transforms of the kernels}
\label{p:app:transforms}

Write $s:=r/\theta$. Both transforms reduce, under $v=\ee^{-\theta t}$, to
integrals of the form $\int_0^1 v^{s-1}(1-v)^n(A+Bv)^{-m}\dd v$, which are
hypergeometric. For the survival kernel, the partial-fraction decomposition
$(a-b)/[(A+Bv)(a-bv)]=B/(A+Bv)+b/(a-bv)$ and the standard integral
$\int_0^1 v^s/(1-zv)\dd v=(s+1)^{-1}\Fhyp(1,s+1;s+2;z)$ give
\begin{equation}
\Lap{\Ifix}(r)=\frac{1}{\lambda(s+1)}
\Biggl[
\Psi(A,B,s)
+\frac{b}{a}\,\Fhyp\Bigl(1,s+1;s+2;\frac{b}{a}\Bigr)
\Biggr],
\label{p:eq:LIhat}
\end{equation}
where, so that the hypergeometric series always converges,
\begin{equation}
\Psi(A,B,s)=
\begin{cases}
\Fhyp\bigl(1,s+1;s+2;-\frac{A}{B}\bigr), & B\ge A,\\[0.4em]
\frac{B}{A}\,\Fhyp\bigl(1,s+1;s+2;-\frac{B}{A}\bigr), & A>B,
\end{cases}
\label{p:eq:Psi}
\end{equation}
the two branches being equal because both sides equal
$(s+1)B\int_0^1 v^s/(A+Bv)\dd v$. For the release kernel the same
substitution and termwise integration give
\begin{equation}
\delta\Lap{K}(r)=\frac{\delta}{\lambda}
\Biggl[
\frac{\Fhyp(1,s;s+2;z)}{A(s+1)}
+\frac{2\,\Fhyp(2,s;s+3;z)}{A^2(s+1)(s+2)}
\Biggr],
\qquad z=-\frac{B}{A},
\label{p:eq:dLK}
\end{equation}
again with the understanding that for $B>A$ the right-hand side is evaluated
through its analytic continuation, or through the integral representation
from which it came. Both reduce correctly at $s=0$: with
$\Fhyp(1,1;2;z)=-\log(1-z)/z$, \cref{p:eq:LIhat} gives
$\Lap{\Ifix}(0)=\lambda^{-1}\log(a/(a-1))=\mean{T_{\rm prod}}$, and
\cref{p:eq:dLK} gives $\delta\Lap{K}(0)=(\delta/\lambda)(1/A+1/A^2)=aB/A
=V_\infty$, using $AB=\delta/\lambda$. Both reductions were verified against
direct quadrature to relative error $\sim10^{-8}$. Where the transforms are
needed as functions of time, standard numerical inversions apply
\cite{stehfest1970algorithm,gaver1965observations}. The curves of
\cref{p:fig:peff} are computed by direct quadrature.

\section{Master formula table}
\label{p:app:formulae}

Throughout: $w=\ee^{\theta t}$, $\theta=\lambda(a-b)$, $A=a-1$, $B=1-b$,
$L=a(1-b)$, $\kappa=1+\delta/(2\lambda)$.

\begin{table}[htbp]
\centering
\renewcommand{\arraystretch}{1.1}
\small
\begin{tabular}{@{}p{0.24\textwidth}p{0.62\textwidth}@{}}
\toprule
Quantity & Formula \\
\midrule
Renewal system & $\Icell=\Icell_0\Ifix+i*\Ifix$, \ $\Vfree'=\Icell_0 g+i*g-c\Vfree$, \ $i=\gamma T\Vfree$ \\
Survival kernel & $S(\alpha)=\Ifix(\alpha)$ \\
Release kernel & $g(\alpha)=\delta K(\alpha)=V'(\alpha)$ \\
Generation kernel & $\mathcal A(\alpha)=\gamma T(g*\ee^{-c\,\cdot})(\alpha)$, \ $\Lap{\mathcal A}(r)=\gamma T\Lap{g}(r)/(r+c)$ \\
$p_{\rm eff}(r)$ & $\delta\Lap{K}(r)/\Lap{\Ifix}(r)$ \\
$d_{\Icell,{\rm eff}}(r)$ & $1/\Lap{\Ifix}(r)-r$ \\
Limits $r=0$ & $p_{\rm eff}=\lambda V_\infty/\log(a/(a-1))$, \ $d_{\rm eff}=\lambda/\log(a/(a-1))$ \\
Limits $r\to\infty$ & $p_{\rm eff}\to\delta$, \ $d_{\rm eff}\to\mu+\delta$ \\
Limit $r\downarrow-\theta$ & $p_{\rm eff}\to\delta\mean{X^2}_{\rm QS}=\delta a(a+1)/(a-1)^2$ \\
Span of the map & $p_{\rm eff}(-\theta)/p_{\rm eff}(\infty)=\mean{X^2}_{\rm QS}$ \\
Characteristic eq. & $r+c=\gamma T\,\delta\Lap{K}(r)$, \ equivalently $\Lap{\mathcal A}(r)=1$ \\
$R_0$ & $\gamma T V_\infty/c=\int_0^\infty\mathcal A$ \\
Branching mean & $m=qV_\infty$, \ $q=\gamma T/(\gamma T+c)$ \\
$G_{\rm off}(z)$ & $b+(\delta/\lambda)\,y/(a-y)$, \ $y=1-q+qz$ \\
$z_{\rm ext}^{\rm burst}$ & $[a(1-q+qb)-1+q]/q$ \\
$z_{\rm ext}^{\rm bud}$ & $1/m$ (geometric output, matched mean) \\
Flooding parameter & $L=a(1-b)=a-\mu/\lambda$, \ $V_\infty=L/(a-1)$ \\
Flooding criterion & $1/\delta<1/\lambda+1/\mu$ \ ($L>1$) \\
$\Delta z_{\rm ext}$ & $(L-1)(1/m-1)$ \\
$\mathrm{Var}_{\rm bud}-\mathrm{Var}_{\rm burst}$ & $2q^2V_\infty(L-1)/(a-1)$ \\
$L=1$ & offspring law exactly geometric; variances equal \\
Growth trade-off & $r_{\rm bud}>r_{\rm burst}$ at matched $R_0$, $\mean{T_{\rm prod}}$ \\
\bottomrule
\end{tabular}
\caption{Population-level quantities, in one place.}
\label{p:tab:formulae}
\end{table}

\paragraph{Specialisation $\mu=0$.} $b=0$, $a=1+\delta/\lambda$, $D\equiv0$,
$\Ifix=I$, $V_\infty=1+\lambda/\delta=\mean{X}_{\rm QS}$,
$\mean{T_{\rm prod}}=\lambda^{-1}\log(1+\lambda/\delta)$,
$\pr{\Kb=k}=s\rho^{k-1}$ with $s=\delta/(\lambda+\delta)$, $\rho=1-s$;
flooding favours bursting for establishment at every $\delta>0$.

\paragraph{The boundary $L=1$.} $\delta=\lambda\mu/(\lambda+\mu)$,
$b=\delta/\lambda$; the offspring law is geometric with mean $m$;
$z_{\rm ext}^{\rm burst}=z_{\rm ext}^{\rm bud}$ for every $q$; offspring
variances equal.

\endgroup
